% !TEX root = ../main.tex
% Chapter 2: Hardy spaces and Nevanlinna class.
% Included from main.tex via % !TEX root = ../main.tex
% Chapter 2: Hardy spaces and Nevanlinna class.
% Included from main.tex via % !TEX root = ../main.tex
% Chapter 2: Hardy spaces and Nevanlinna class.
% Included from main.tex via % !TEX root = ../main.tex
% Chapter 2: Hardy spaces and Nevanlinna class.
% Included from main.tex via \input{chapters/Chapter2}.

\chapter{Hardy spaces and Nevanlinna class}
\label{sec:hardy_nevanlinna}
This chapter gives an introduction on the mathematical backgrounds of
the main objects of this thesis: Hardy spaces and 
Nevanlinna class, and discuss the ultimate goal to formalize in
this project: canonical factorization theorem of Nevanlinna functions. This chapter is
mainly based on \cite{garnett2006bounded}, II.1 - II.5.

\section{Hardy spaces}
One of the main objects of this thesis is the \textbf{Hardy space}, which is a subspace of the space 
$\mathcal{O}(U)$ of analytic functions
on a `certain domain' $U\subset \mathbb{C}$ which are defined by 
`certain growth conditions'. Roughly speaking,
`certain growth conditions' means that the boundary growth is controlled in 
an averaged $L^p$ sense. We usually consider
$U$ to be the unit disc $\D:=\{z\in\C: |z|<1\}$ or the upper half plane 
$\HH:=\{z\in\C: \im z>0\}$.

\begin{definition}
Let $0<p<\infty$. We define the \textbf{Hardy space over the unit disc $\D$} as:
\begin{equation*}
  H^p(\D):=\left\{ f\in \mathcal{O}(\D):\| f\|_{H^p}^p := \sup_{0<r<1}\frac{1}{2\pi}\int |f(re^{i\theta})|^p d
  \theta <\infty\right\}
\end{equation*}
If $p=\infty$, we define $H^\infty(\D)$ to be the space of bounded analytic functions on $\D$, i.e.
\begin{equation*}
  H^\infty(\D):=\{f\in \mathcal{O}(\D): \|f\|_{H^\infty}:=\sup_{z\in \D}|f(z)|<\infty\}
\end{equation*}
\end{definition}

\begin{definition}
Let $0<p<\infty$. We define the \textbf{Hardy space over the upper half plane $\HH$} as:
\begin{equation*}
  H^p(dt):=\left\{f\in \mathcal{O}(\HH): \|f\|_{H^p}^p:=\sup_{y>0}\int |f(x+iy)|^p dx <\infty\right\}
\end{equation*}
If $p=\infty$, we define $H^\infty(dt)$ to be the space of bounded analytic functions on $\HH$, i.e.
\begin{equation*}
  H^\infty(dt):=\{f\in \mathcal{O}(\HH): \|f\|_{H^\infty}:=\sup_{z\in \HH}|f(z)|<\infty\}
\end{equation*}
\end{definition}
These two definitions are closely related and many results holds true on $H^p(\D)$ will be also holds
true on $H^p(dt)$. In this thesis, we only focus on the formalization for $H^p(\D)$.

\begin{remark}
One can extend the definition of Hardy space over general (bounded) domain in $\mathbb{C}$ as follow:
\begin{definition}
Let $0<p<\infty$ and $U\subset \mathbb{C}$ be a (bounded) domain. We say an analytic function $f\in \mathcal{O}(U)$
over $U$ is a $H^p$ function if the function $|f|^p$ has a harmonic majorant, and we define $\|f\|^p_{H^p}$
to be the value of the least harmonic majorant at $z=0$. For $p=\infty$, $f\in \mathcal{O}(U)$ is
$H^\infty$ if it is a bounded analytic function over $U$ and $\|f\|_{H^\infty}:=\sup_{z\in U}|f(z)|$.
\end{definition}
Theorem \ref{Harmonic majorant criterion} shows that this is an equivalent definition with above definitions over $\D$.
\end{remark}


The Hardy space $H^p(\D)$ is a normed space for $p\geq 1$ with the
norm $\|\cdot\|_{H^p}$. Moreover, it is a metric space for all $0<p\leq \infty$,
with the metric $d(f,g)=\|f-g\|_{H^p}^p$ for $0<p<1$ and the metric induced by
the norm for $p\geq 1$. Moreover, we have:
\begin{theorem}
For $0<p\leq \infty$, $H^p$ is a complete metric space.
\end{theorem}
Similar as $L^p$ space, there is monotonicity for $H^p$ space as well:
\begin{theorem}
For $0<p\leq q\leq \infty$, we have $H^q\leq H^p$ as subspace.
\end{theorem}


\section{Blaschke product}
\begin{definition}
Given a sequence $\{z_n\}$ on $\mathbb{D}$, we say that this sequence satisfies
\textbf{Blaschke's condition} if $\sum (1-|z_n|)<\infty$. Given an analytic function $f$ on 
$\mathbb{D}$, we call this function $f$ satisfies \textbf{Blaschke's condition} if its
zeros $\{z_n\}$ satisfies Blaschke's condition.
\end{definition}

\begin{theorem}\label{Blaschke condition from harmonic majorant}
Let $f\not\equiv 0$ be an analytic function on $\mathbb{D}$ with zeros $\{z_n\}$. If
$\log|f|$ has a harmonic majorant, then $f$ satisfies Blaschke's condition:
\begin{equation*}
  \sum(1-|z_n|)<\infty.
\end{equation*}
If moreover $f(0)\neq 0$ and $u$ is the least harmonic majorant of $\log|f|$, then
\begin{equation*}
  \sum(1-|z_n|)\leq u(0)-\log|f(0)|.
\end{equation*}
\end{theorem}

\begin{theorem}\label{Blaschke product}
Let $\{z_n\}$ be a sequence in $\mathbb{D}$ satisfying Blaschke's condition, and let $m$ be the
number of $z_n$ equal to $0$. Then the \textbf{Blaschke product}
\begin{equation}\label{eq:Blaschke product}
  B(z)=z^m\prod_{|z_n|\neq 0}\left(-\frac{\overline{z_n}}{|z_n|}\right)
  \frac{z-z_n}{1-\overline{z_n}z}
\end{equation}
converges on $\mathbb{D}$. The function $B$ is in $H^\infty(\mathbb{D})$ with $|B(z)|\leq 1$,
and the zeros of $B$ are precisely the points $z_n$, each zero having multiplicity equal to
the number of times it occurs in $\{z_n\}$. Moreover,
\begin{equation*}
  |B(e^{i\theta})|=1
\end{equation*}
for almost every $\theta$, where $B(e^{i\theta})=\lim_{r\rightarrow 1}B(re^{i\theta})$ is the
radial limit of $B$.
\end{theorem}

Here's a characterization of Blaschke products in terms of harmonic majorants:
\begin{theorem}\label{characterization of Blaschke products}
Let $f\in H^\infty(\mathbb{D})$ with $\|f\|_{H^\infty}\leq 1$. The followings are equivalent:
\begin{itemize}
  \item[(a)] $f(z)=\lambda B(z)$ where $\lambda$ is constant with $|\lambda|=1$ and $B$ is a 
  Blaschke product;
  \item[(b)] $\lim_{r\rightarrow 1}\int \log |f(re^{i\theta})|d\theta/(2\pi)=0$;
  \item[(c)] The least harmonic majorant of $\log |f|$ is 0.
\end{itemize}
\end{theorem}

\section{Boundary functions of Hardy functions}
Every $H^p$ function has a boundary function in $L^p(\mathbb{T})$, and for $p\geq 1$ it can be
recovered from its boundary function via Poisson integral.
\begin{theorem}\label{thm:Hp_boundary_function}
Let $0<p\leq\infty$ and $f\in H^p(\D)$. Then $f$ has a nontangential limit $f^*(\zeta)$ at
almost every $\zeta\in\mathbb{T}$, and $f^*\in L^p(\mathbb{T}, \frac{d\theta}{2\pi})$ with
\begin{equation*}
  \|f^*\|_{L^p}=\|f\|_{H^p}.
\end{equation*}
Moreover:
\begin{itemize}
  \item[(a)] if $0<p<\infty$, then
  \begin{equation*}
    \lim_{r\rightarrow 1}\|f(re^{i\cdot})-f^*\|_{L^p}=0;
  \end{equation*}
  \item[(b)] if $1\leq p\leq\infty$, then $f=P[f^*]$.
\end{itemize}
\end{theorem}




\section{Nevanlinna class} 
\begin{definition}
An analytic function $f(z)$ on $\D$ is in the \textbf{Nevanlinna class} $N$ if 
\begin{equation*}
  \sup_r \int \log^+|f(r e^{i\theta})|\frac{d\theta}{2\pi}<\infty
\end{equation*}
\end{definition}
As $\log^+|f(z)|\leq (1/p)|f(z)|^p$, so $H^p\subset N$ for all $0<p\leq \infty$, i.e. Nevanlinna
class encompasses all Hardy spaces $H^p$. 

By Theorem \ref{Harmonic majorant criterion}, we can formulate the definition
in an equivalent definition:
\begin{definition}
An analytic function $f(z)$ on $\D$ is in the \textbf{Nevanlinna class} $N$ if 
the subharmonic function $\log^+|f(z)|$ has a harmonic majorant.
\end{definition}

\begin{remark}\label{rem:N_algebra}
Compared to $H^p(\D) \subset \mathcal{O}(\D)$,
which is just a subspace of $\mathcal{O}(\D)$ for $p<\infty$, $N\subset \mathcal{O}(\D)$ is
also a subalgebra of $\mathcal{O}(\D)$: if $f,g\in N$, then $fg\in N$ as
\begin{equation*}
  \log^+|fg|\leq \log^+|f|+\log^+|g|.
\end{equation*}
Moreover, if $f\in N$ has no zeros in $\D$, then $1/f\in N$. Indeed, let $U$ be a harmonic
majorant of $\log^+|f|$. Since $f$ has no zeros, $\log|f|$ is harmonic on $\D$, and
\begin{equation*}
  \log^+|1/f|=\log^+|f|-\log|f|\leq U-\log|f|,
\end{equation*}
so $\log^+|1/f|$ has the harmonic majorant $U-\log|f|$.
The closure of multiplication and inversion will sort of provide a more natural place for us to 
describe `factorization problem'.
\end{remark}

\begin{theorem}\label{characterization of Nevanlinna using log |f|}
Let $f(z)$ be an analytic function on $\mathbb{D}$ s.t. $f\not\equiv 0$. 
Then $f\in N$ if and only if $\log|f(z)|$ has least harmonic majorant
the Poisson integral of some finite measure on $\mathbb{T}$.
\end{theorem}
\begin{proof}

\end{proof}


\begin{lemma}\label{least harmonic majorant of log|f| for Nevanlinna f}
Let $f\in N$ and $f\not\equiv 0$. Let $B$ be the Blaschke product formed from zeros of $f$. Then:
\begin{itemize}
  \item $B$ converges;
  \item $g=f/B\in N$;
  \item $\log |g|$ is the least harmonic majorant of $\log |f|$.
\end{itemize}
\end{lemma}
\begin{proof}[Sketch of proof]
As $f\in N$, $\log |f|$ has a harmonic majorant, thus by Theorem \ref{Blaschke condition from harmonic majorant},
$f$ satisfies Blaschke's condition, thus $B$ converges. Now let $u$ be the least harmonic majorant
of $\log |f|$. Recall in Example \ref{log |f| is subharmonic}, $g$ as a zero-free analytic function
over $\mathbb{D}$, $\log |g|$ is harmonic on $\mathbb{D}$. As $|B(z)|\leq 1$, we have 
$\log |f(z)|\leq \log |g(z)|$, thus 
\begin{equation*}
  u(z)\leq \log|g(z)|.
\end{equation*}
On the other hand, 
\begin{equation*}
  \log |B(z)|=\log |f(z)|-\log|g(z)|\leq u(z)-\log |g(z)|
\end{equation*}
Recall Theorem \ref{characterization of Blaschke products}, the least harmonic majorant of $\log|B|$
is 0, thus 
\begin{equation*}
  0\leq u(z)-\log|g(z)|
\end{equation*}
which implies $u(z)=\log |g(z)|$. It's easy to check that $g\in N$ by definition.
\end{proof}


\begin{theorem}\label{nontangential limit of Nevanlinna}
Let $f\in N$ and $f\not\equiv 0$. Then $f$ has a nontangential limit $f^*(\zeta)$ at almost
every $\zeta\in\mathbb{T}$, and $\log|f^*|\in L^1(\mathbb{T})$.

Moreover, the least harmonic majorant of $\log|f|$ is $P[\mu]$, where
\begin{equation*}
  d\mu=\log|f^*|\,\frac{d\theta}{2\pi}+d\mu_s,
  \qquad \mu_s\perp\frac{d\theta}{2\pi}.
\end{equation*}
\end{theorem}

\subsection{Outer functions and Singular functions}
Before we discuss the main topic of this section, which is the factorization problem of
Nevanlinna functions, we'd like to first introduce two classes of functions: 
\textbf{outer functions} and \textbf{singular functions}.
\begin{definition}\label{def:outer_function}
An analytic function $F$ on $\D$ is an \textbf{outer function} if
\begin{equation*}
  F(z)=C\exp\left(\int K_z(\theta)\,l(\theta)\,\frac{d\theta}{2\pi}\right),
  \qquad z\in\D,
\end{equation*}
for some real-valued $l\in L^1(\mathbb{T})$ and some constant $C\in\mathbb{C}$ with $|C|=1$.
Here $K_z(\theta)=\frac{e^{i\theta}+z}{e^{i\theta}-z}$ is the Herglotz-Riesz kernel.
\end{definition}

\begin{theorem}\label{thm:outer_function_properties}
Let $F$ be an outer function on $\mathbb{D}$. Then:
\begin{itemize}
  \item[(a)] $F$ is analytic on $\D$ and has no zeros in $\D$;
  \item[(b)] $\log|F|=P[l]$ on $\D$;
  \item[(c)] $F\in N$;
  \item[(d)] $F$ has a nontangential limit $F^*(\zeta)$ at almost every $\zeta\in\mathbb{T}$, and
  $|F^*|=e^{l}$ almost everywhere on $\mathbb{T}$;
  \item[(e)] if $F_1,F_2$ are outer functions with $|F_1^*|=|F_2^*|$ almost everywhere on
  $\mathbb{T}$, then $F_1=cF_2$ for some constant $c\in\mathbb{C}$ with $|c|=1$.
\end{itemize}
\end{theorem}
\begin{proof}
Write $F=Ce^{G}$ with
\begin{equation*}
  G(z)=\int K_z(\theta)\,l(\theta)\,\frac{d\theta}{2\pi}.
\end{equation*}

(a) For $|z|\leq r<1$ we have $|K_z(\theta)|\leq\frac{1+r}{1-r}$, and $K_z(\theta)$ is analytic
in $z$. Hence $G$ is analytic on $\D$ and so is
$F=Ce^{G}$. Since the exponential has no zeros, neither does $F$.

(b) As $|C|=1$, $l$ is real-valued and $\text{Re } K_z(\theta)=P_z(\theta)$, note
that $\log |e^G|=\text{Re } G$,

\begin{equation*}
  \log|F(z)|=\text{Re } G(z)=\int P_z(\theta)\,l(\theta)\,\frac{d\theta}{2\pi}=P[l](z).
\end{equation*}

(c) Write $l^+=\max(l,0)$. By Lemma \ref{Properties of Poisson kernel}(a), $P_z>0$, so
$\log|F|=P[l]\leq P[l^+]$, and since $P[l^+]\geq 0$, also $\log^+|F|\leq P[l^+]$. Using
$P_{re^{i\theta}}(\varphi)=P_{re^{i\varphi}}(\theta)$, Fubini's theorem and Lemma
\ref{Properties of Poisson kernel}(a) again,
\begin{equation*}
  \frac{1}{2\pi}\int_0^{2\pi}\log^+|F(re^{i\theta})|\,d\theta
  \leq\frac{1}{2\pi}\int_0^{2\pi}P[l^+](re^{i\theta})\,d\theta
  =\frac{1}{2\pi}\int_0^{2\pi}l^+(\varphi)\,d\varphi
  \leq\|l\|_{L^1}
\end{equation*}
for every $0<r<1$. Hence $F\in N$.

(d) By (c) and $F\not\equiv0$, Theorem \ref{nontangential limit of Nevanlinna} shows that $F$ has
a nontangential limit $F^*$ almost everywhere. By (b) and Corollary
\ref{cor:poisson_nontangential_ae}, $\log|F|=P[l]$ has nontangential limit $l$ almost
everywhere. At almost every $\zeta\in\mathbb{T}$ both limits exist, and then
\begin{equation*}
  |F^*(\zeta)|=\lim_{\substack{z\to\zeta\\ z\in\Gamma_\alpha(\zeta)}}e^{\log|F(z)|}=e^{l(\zeta)}.
\end{equation*}

(e) Let $F_j=C_je^{G_j}$, where $G_j$ is defined from $l_j$ as above, $j=1,2$. By (d),
$e^{l_1}=|F_1^*|=|F_2^*|=e^{l_2}$ almost everywhere, so $l_1=l_2$ almost everywhere. Hence
$G_1=G_2$ on $\D$, and $F_1=cF_2$ with $c=C_1/C_2$, $|c|=1$.
\end{proof}

\begin{remark}
One can construct an outer function as follow: given $h:\mathbb{T}\rightarrow [0,\infty)$
s.t. $\log h\in L^1(\mathbb{T})$
(think of it as the modulos of the boundary function of Nevanlinna functions over $\mathbb{D}$, 
see \ref{nontangential limit of Nevanlinna}), then set $u(z)=P[\log h](z)$ and $v(z)$ to be one of the
harmonic conjugates of $u(z)$. Then the function $F(z):= e^{u(z)+i v(z)}$ is an outer function,
and it is determined by $h(\theta)$ except for a unimodular constant factor due to the choice
of the conjugate $v(z)$. 
\end{remark}

\begin{definition}
An analytic function $S$ on $\D$ is a \textbf{singular function} if
\begin{equation*}
  S(z)=\exp\left(-\int K_z(\theta)\,d\mu(\theta)\right),
  \qquad z\in\D,
\end{equation*}
for some finite singular positive measure $\mu$ on $\mathbb{T}$.
\end{definition}

\begin{lemma}\label{lem:log_singular_function}
For singular function $S(z)=\exp\left(-\int K_z(\theta)\,d\mu(\theta)\right)$
 on $\mathbb{D}$, we have $\log|S|=-P[\mu]$ on $\mathbb{D}$.
\end{lemma}
\begin{proof}
Given $z\in\D$, as $|K_z(\theta)|\leq\frac{1+|z|}{1-|z|}$ and $\mu$ is finite, 
$\int K_z(\theta)\,d\mu(\theta)$ converges. As $\log|e^w|=\text{Re } w$, $\mu$ is a positive
measure and $\text{Re } K_z(\theta)=P_z(\theta)$,
\begin{equation*}
  \log|S(z)|=-\text{Re }\int K_z(\theta)\,d\mu(\theta)
  =-\int P_z(\theta)\,d\mu(\theta)=-P[\mu](z).
\end{equation*}
\end{proof}

\begin{lemma}\label{lem:singular_function_in_N}
If $S$ is a singular function on $\mathbb{D}$, then $S\in N$.
\end{lemma}
\begin{proof}
By Lemma \ref{lem:log_singular_function} and Lemma \ref{Properties of Poisson kernel}(a),
$\log|S|=-P[\mu]\leq 0$ on $\D$, as $\mu\geq 0$ and $P_z>0$. Hence $\log^+|S|=0$ on $\D$,
so $S\in N$.
\end{proof}

\begin{theorem}\label{equivalent characterization of singular function}
Let $S$ be an analytic function on $\D$. Then $S$ is a singular function if and only if $S$ has the
following properties:
\begin{itemize}
  \item[(a)] $S$ has no zeros in $\D$;
  \item[(b)] $|S(z)|\leq 1$ for all $z\in\D$;
  \item[(c)] $S$ has a nontangential limit $S^*(\zeta)$ at almost every $\zeta\in\mathbb{T}$, and
  $|S^*|=1$ almost everywhere on $\mathbb{T}$;
  \item[(d)] $S(0)>0$.
\end{itemize}
\end{theorem}

\begin{proof}[Proof of one direction]
($\Rightarrow$) Let $S(z)=\exp\left(-\int K_z(\theta)\,d\mu(\theta)\right)$ with $\mu\geq 0$
finite and singular.
(a) holds as exponential has no zeros.
(b) was shown in the proof of Lemma \ref{lem:singular_function_in_N}.
(c) By Lemma \ref{lem:singular_function_in_N} and Theorem \ref{nontangential limit of Nevanlinna},
$S$ has a nontangential limit $S^*$ almost everywhere. By Lemma \ref{lem:log_singular_function}
and Theorem \ref{thm:singular_poisson_nontangential_zero}, $\log|S|=-P[\mu]$ has nontangential
limit $0$ almost everywhere. So at a.e. $\zeta\in\mathbb{T}$,
\begin{equation*}
  |S^*(\zeta)|=\lim_{\substack{z\to\zeta\\ z\in\Gamma_\alpha(\zeta)}}e^{-P[\mu](z)}=1.
\end{equation*}
(d) holds because $K_0(\theta)=1$ gives $S(0)=e^{-\mu(\mathbb{T})}>0$.

We will skip the proof of ($\Leftarrow$) as it's not needed in the project.
\end{proof}



\subsection{Canonical factorization theorem of Nevanlinna funcitons}
Now we'd like to discuss the \textit{factorization problem} of Nevanlinna functions.
Given $f\in N$ with $f\not\equiv 0$, consider $B$: Blaschke product formed by
its zeros and $g=f/B$. By Lemma \ref{least harmonic majorant of log|f| for Nevanlinna f},
we know that $\log|g|$ is the least harmonic majorant of $\log |f|$. Recall 
Theorem \ref{characterization of Nevanlinna using log |f|}, we can write 
$\log |g|=P[\mu]$ for some finite measure $\mu$. By Theorem \ref{nontangential limit of Nevanlinna},
we can write:
\begin{equation*}
  d\mu=\log|f^*|\,\frac{d\theta}{2\pi}+d\mu_s,
  \qquad \mu_s\perp\frac{d\theta}{2\pi}.
\end{equation*}
Now take Jordan decomposition on singular measure
 $\mu_s$, we get the unique decomposition:
 \begin{equation*}
  \mu_s=-\mu_1+\mu_2
 \end{equation*}
 where $\mu_1$ and $\mu_2$ are positive singular measure that are mutually singular to each other. Thus we 
write:
\begin{equation*}
  d\mu=\log|f^*|\,\frac{d\theta}{2\pi}-d\mu_1+d\mu_2.
\end{equation*}

Now consider the \textbf{outer function}
\begin{equation*}
  \tilde{F}(z)=\exp\left(\int K_z(\theta)\log |f^*(e^{i\theta})|\frac{d\theta}{2\pi}\right)
\end{equation*}
Here $K_z(\theta)=\frac{e^{i\theta}+z}{e^{i\theta}-z}$ is the Herglotz-Riesz kernel. Moreover, consider
the \textbf{singular functions}:
\begin{equation*}
  S_j(z)=\exp\left(-\int K_z(\theta)d\mu_j(\theta)\right),\quad j=1,2.
\end{equation*}
Recall that $\text{Re } K_z(\theta)=P_z(\theta)$, which is the Poisson kernel, and
$\log |e^w|=\text{Re } w$, so 
\begin{equation*}
  \log |\tilde{F}(z)|=P[\log |f|](z),\; \log |S_j(z)|=- P[\mu_j](z),\quad j=1,2.
\end{equation*}
Therefore, 
\begin{align*}
  \log |g(z)|=P[\mu](z)&=P[\log |f|](z)+P[\mu_1](z)-P[\mu_2](z)\\
  &= \log |\tilde{F}(z)|+\log |S_1(z)|-\log|S_2(z)|.
\end{align*}
Then we have:
\begin{equation*}
  \log \left|\frac{g(z) S_2(z)}{\tilde{F}(z)S_1(z)}\right|=0
\end{equation*}
which takes $\exp$ on both sides yields:
\begin{equation*}
  \left|\frac{f(z)S_2(z)}{B(z)\tilde{F}(z)S_1(z)}\right|=1,\quad \forall z\in \mathbb{D}.
\end{equation*}
By maximum modulus principle, 
\begin{equation*}
  \frac{f(z)S_2(z)}{B(z)\tilde{F}(z)S_1(z)}=C,\quad C\in\mathbb{C},\; |C|=1.
\end{equation*}
Set $F=C\tilde{F}$. By Definition \ref{def:outer_function}, $F$ is again an outer function, and
$f=BFS_1/S_2$.
This gives most of the \textbf{canonical factorization theorem}:
\begin{theorem}
Let $f\in N$, $f\not\equiv 0$. Then
\begin{equation*}
  f(z)=B(z)F(z)S_1(z)/S_2(z),
\end{equation*}
where $B$ is Blaschke product formed by zeros of $f$,
$F$ is an outer function, and $S_1,S_2$ are
singular functions whose positive singular measures $\mu_1,\mu_2$ are mutually singular.
This factorization is unique.
Conversely, every function of this form is in $N$.
\end{theorem}
\begin{proof}
The existence part is given as above.

For uniqueness, the Blaschke product $B$ is determined by the zeros of $f$. Suppose
$f=BFS_1/S_2=BF'S_1'/S_2'$, where $S_1,S_2$ come from $\mu_1,\mu_2$ and $S_1',S_2'$ from
$\nu_1,\nu_2$. Since the zeros of $B\not\equiv 0$ are isolated, $FS_1/S_2=F'S_1'/S_2'$ on $\D$
minus a discrete set, hence on $\D$ by continuity. By Theorem
\ref{equivalent characterization of singular function}(c), $|S_1^*|=|S_2^*|=|S_1'^*|=|S_2'^*|=1$
almost everywhere, so $|F^*|=|F'^*|$ almost everywhere, and by Theorem
\ref{thm:outer_function_properties}(e), $F=cF'$ for a constant $c$ with $|c|=1$. Thus
$cS_1/S_2=S_1'/S_2'$, and taking $\log|\cdot|$, Lemma \ref{lem:log_singular_function} gives
\begin{equation*}
  P[\mu_2-\mu_1]=P[\nu_2-\nu_1]\quad\text{on }\D.
\end{equation*}
Since a finite measure on $\mathbb{T}$ is determined by its Poisson integral,
$\mu_2-\mu_1=\nu_2-\nu_1$. Both $(\mu_2,\mu_1)$ and $(\nu_2,\nu_1)$ are pairs of mutually
singular positive measures, so by uniqueness of the Jordan decomposition, $\mu_1=\nu_1$ and
$\mu_2=\nu_2$. Hence $S_j=S_j'$ for $j=1,2$, so $c=1$ and $F=F'$.

For converse part Since $N$ is closed under multiplication by Remark \ref{rem:N_algebra}, it suffices to show that
$B,F,S_1,1/S_2\in N$. We have $B\in H^\infty\subset N$ by Theorem \ref{Blaschke product},
$F\in N$ by Theorem \ref{thm:outer_function_properties}(c), and $S_1\in N$ by Lemma
\ref{lem:singular_function_in_N}. Likewise $S_2\in N$, and $S_2$ has no zeros in $\D$, so
$1/S_2\in N$ by Remark \ref{rem:N_algebra}.
\end{proof}
.

\chapter{Hardy spaces and Nevanlinna class}
\label{sec:hardy_nevanlinna}
This chapter gives an introduction on the mathematical backgrounds of
the main objects of this thesis: Hardy spaces and 
Nevanlinna class, and discuss the ultimate goal to formalize in
this project: canonical factorization theorem of Nevanlinna functions. This chapter is
mainly based on \cite{garnett2006bounded}, II.1 - II.5.

\section{Hardy spaces}
One of the main objects of this thesis is the \textbf{Hardy space}, which is a subspace of the space 
$\mathcal{O}(U)$ of analytic functions
on a `certain domain' $U\subset \mathbb{C}$ which are defined by 
`certain growth conditions'. Roughly speaking,
`certain growth conditions' means that the boundary growth is controlled in 
an averaged $L^p$ sense. We usually consider
$U$ to be the unit disc $\D:=\{z\in\C: |z|<1\}$ or the upper half plane 
$\HH:=\{z\in\C: \im z>0\}$.

\begin{definition}
Let $0<p<\infty$. We define the \textbf{Hardy space over the unit disc $\D$} as:
\begin{equation*}
  H^p(\D):=\left\{ f\in \mathcal{O}(\D):\| f\|_{H^p}^p := \sup_{0<r<1}\frac{1}{2\pi}\int |f(re^{i\theta})|^p d
  \theta <\infty\right\}
\end{equation*}
If $p=\infty$, we define $H^\infty(\D)$ to be the space of bounded analytic functions on $\D$, i.e.
\begin{equation*}
  H^\infty(\D):=\{f\in \mathcal{O}(\D): \|f\|_{H^\infty}:=\sup_{z\in \D}|f(z)|<\infty\}
\end{equation*}
\end{definition}

\begin{definition}
Let $0<p<\infty$. We define the \textbf{Hardy space over the upper half plane $\HH$} as:
\begin{equation*}
  H^p(dt):=\left\{f\in \mathcal{O}(\HH): \|f\|_{H^p}^p:=\sup_{y>0}\int |f(x+iy)|^p dx <\infty\right\}
\end{equation*}
If $p=\infty$, we define $H^\infty(dt)$ to be the space of bounded analytic functions on $\HH$, i.e.
\begin{equation*}
  H^\infty(dt):=\{f\in \mathcal{O}(\HH): \|f\|_{H^\infty}:=\sup_{z\in \HH}|f(z)|<\infty\}
\end{equation*}
\end{definition}
These two definitions are closely related and many results holds true on $H^p(\D)$ will be also holds
true on $H^p(dt)$. In this thesis, we only focus on the formalization for $H^p(\D)$.

\begin{remark}
One can extend the definition of Hardy space over general (bounded) domain in $\mathbb{C}$ as follow:
\begin{definition}
Let $0<p<\infty$ and $U\subset \mathbb{C}$ be a (bounded) domain. We say an analytic function $f\in \mathcal{O}(U)$
over $U$ is a $H^p$ function if the function $|f|^p$ has a harmonic majorant, and we define $\|f\|^p_{H^p}$
to be the value of the least harmonic majorant at $z=0$. For $p=\infty$, $f\in \mathcal{O}(U)$ is
$H^\infty$ if it is a bounded analytic function over $U$ and $\|f\|_{H^\infty}:=\sup_{z\in U}|f(z)|$.
\end{definition}
Theorem \ref{Harmonic majorant criterion} shows that this is an equivalent definition with above definitions over $\D$.
\end{remark}


The Hardy space $H^p(\D)$ is a normed space for $p\geq 1$ with the
norm $\|\cdot\|_{H^p}$. Moreover, it is a metric space for all $0<p\leq \infty$,
with the metric $d(f,g)=\|f-g\|_{H^p}^p$ for $0<p<1$ and the metric induced by
the norm for $p\geq 1$. Moreover, we have:
\begin{theorem}
For $0<p\leq \infty$, $H^p$ is a complete metric space.
\end{theorem}
Similar as $L^p$ space, there is monotonicity for $H^p$ space as well:
\begin{theorem}
For $0<p\leq q\leq \infty$, we have $H^q\leq H^p$ as subspace.
\end{theorem}


\section{Blaschke product}
\begin{definition}
Given a sequence $\{z_n\}$ on $\mathbb{D}$, we say that this sequence satisfies
\textbf{Blaschke's condition} if $\sum (1-|z_n|)<\infty$. Given an analytic function $f$ on 
$\mathbb{D}$, we call this function $f$ satisfies \textbf{Blaschke's condition} if its
zeros $\{z_n\}$ satisfies Blaschke's condition.
\end{definition}

\begin{theorem}\label{Blaschke condition from harmonic majorant}
Let $f\not\equiv 0$ be an analytic function on $\mathbb{D}$ with zeros $\{z_n\}$. If
$\log|f|$ has a harmonic majorant, then $f$ satisfies Blaschke's condition:
\begin{equation*}
  \sum(1-|z_n|)<\infty.
\end{equation*}
If moreover $f(0)\neq 0$ and $u$ is the least harmonic majorant of $\log|f|$, then
\begin{equation*}
  \sum(1-|z_n|)\leq u(0)-\log|f(0)|.
\end{equation*}
\end{theorem}

\begin{theorem}\label{Blaschke product}
Let $\{z_n\}$ be a sequence in $\mathbb{D}$ satisfying Blaschke's condition, and let $m$ be the
number of $z_n$ equal to $0$. Then the \textbf{Blaschke product}
\begin{equation}\label{eq:Blaschke product}
  B(z)=z^m\prod_{|z_n|\neq 0}\left(-\frac{\overline{z_n}}{|z_n|}\right)
  \frac{z-z_n}{1-\overline{z_n}z}
\end{equation}
converges on $\mathbb{D}$. The function $B$ is in $H^\infty(\mathbb{D})$ with $|B(z)|\leq 1$,
and the zeros of $B$ are precisely the points $z_n$, each zero having multiplicity equal to
the number of times it occurs in $\{z_n\}$. Moreover,
\begin{equation*}
  |B(e^{i\theta})|=1
\end{equation*}
for almost every $\theta$, where $B(e^{i\theta})=\lim_{r\rightarrow 1}B(re^{i\theta})$ is the
radial limit of $B$.
\end{theorem}

Here's a characterization of Blaschke products in terms of harmonic majorants:
\begin{theorem}\label{characterization of Blaschke products}
Let $f\in H^\infty(\mathbb{D})$ with $\|f\|_{H^\infty}\leq 1$. The followings are equivalent:
\begin{itemize}
  \item[(a)] $f(z)=\lambda B(z)$ where $\lambda$ is constant with $|\lambda|=1$ and $B$ is a 
  Blaschke product;
  \item[(b)] $\lim_{r\rightarrow 1}\int \log |f(re^{i\theta})|d\theta/(2\pi)=0$;
  \item[(c)] The least harmonic majorant of $\log |f|$ is 0.
\end{itemize}
\end{theorem}

\section{Boundary functions of Hardy functions}
Every $H^p$ function has a boundary function in $L^p(\mathbb{T})$, and for $p\geq 1$ it can be
recovered from its boundary function via Poisson integral.
\begin{theorem}\label{thm:Hp_boundary_function}
Let $0<p\leq\infty$ and $f\in H^p(\D)$. Then $f$ has a nontangential limit $f^*(\zeta)$ at
almost every $\zeta\in\mathbb{T}$, and $f^*\in L^p(\mathbb{T}, \frac{d\theta}{2\pi})$ with
\begin{equation*}
  \|f^*\|_{L^p}=\|f\|_{H^p}.
\end{equation*}
Moreover:
\begin{itemize}
  \item[(a)] if $0<p<\infty$, then
  \begin{equation*}
    \lim_{r\rightarrow 1}\|f(re^{i\cdot})-f^*\|_{L^p}=0;
  \end{equation*}
  \item[(b)] if $1\leq p\leq\infty$, then $f=P[f^*]$.
\end{itemize}
\end{theorem}




\section{Nevanlinna class} 
\begin{definition}
An analytic function $f(z)$ on $\D$ is in the \textbf{Nevanlinna class} $N$ if 
\begin{equation*}
  \sup_r \int \log^+|f(r e^{i\theta})|\frac{d\theta}{2\pi}<\infty
\end{equation*}
\end{definition}
As $\log^+|f(z)|\leq (1/p)|f(z)|^p$, so $H^p\subset N$ for all $0<p\leq \infty$, i.e. Nevanlinna
class encompasses all Hardy spaces $H^p$. 

By Theorem \ref{Harmonic majorant criterion}, we can formulate the definition
in an equivalent definition:
\begin{definition}
An analytic function $f(z)$ on $\D$ is in the \textbf{Nevanlinna class} $N$ if 
the subharmonic function $\log^+|f(z)|$ has a harmonic majorant.
\end{definition}

\begin{remark}\label{rem:N_algebra}
Compared to $H^p(\D) \subset \mathcal{O}(\D)$,
which is just a subspace of $\mathcal{O}(\D)$ for $p<\infty$, $N\subset \mathcal{O}(\D)$ is
also a subalgebra of $\mathcal{O}(\D)$: if $f,g\in N$, then $fg\in N$ as
\begin{equation*}
  \log^+|fg|\leq \log^+|f|+\log^+|g|.
\end{equation*}
Moreover, if $f\in N$ has no zeros in $\D$, then $1/f\in N$. Indeed, let $U$ be a harmonic
majorant of $\log^+|f|$. Since $f$ has no zeros, $\log|f|$ is harmonic on $\D$, and
\begin{equation*}
  \log^+|1/f|=\log^+|f|-\log|f|\leq U-\log|f|,
\end{equation*}
so $\log^+|1/f|$ has the harmonic majorant $U-\log|f|$.
The closure of multiplication and inversion will sort of provide a more natural place for us to 
describe `factorization problem'.
\end{remark}

\begin{theorem}\label{characterization of Nevanlinna using log |f|}
Let $f(z)$ be an analytic function on $\mathbb{D}$ s.t. $f\not\equiv 0$. 
Then $f\in N$ if and only if $\log|f(z)|$ has least harmonic majorant
the Poisson integral of some finite measure on $\mathbb{T}$.
\end{theorem}
\begin{proof}

\end{proof}


\begin{lemma}\label{least harmonic majorant of log|f| for Nevanlinna f}
Let $f\in N$ and $f\not\equiv 0$. Let $B$ be the Blaschke product formed from zeros of $f$. Then:
\begin{itemize}
  \item $B$ converges;
  \item $g=f/B\in N$;
  \item $\log |g|$ is the least harmonic majorant of $\log |f|$.
\end{itemize}
\end{lemma}
\begin{proof}[Sketch of proof]
As $f\in N$, $\log |f|$ has a harmonic majorant, thus by Theorem \ref{Blaschke condition from harmonic majorant},
$f$ satisfies Blaschke's condition, thus $B$ converges. Now let $u$ be the least harmonic majorant
of $\log |f|$. Recall in Example \ref{log |f| is subharmonic}, $g$ as a zero-free analytic function
over $\mathbb{D}$, $\log |g|$ is harmonic on $\mathbb{D}$. As $|B(z)|\leq 1$, we have 
$\log |f(z)|\leq \log |g(z)|$, thus 
\begin{equation*}
  u(z)\leq \log|g(z)|.
\end{equation*}
On the other hand, 
\begin{equation*}
  \log |B(z)|=\log |f(z)|-\log|g(z)|\leq u(z)-\log |g(z)|
\end{equation*}
Recall Theorem \ref{characterization of Blaschke products}, the least harmonic majorant of $\log|B|$
is 0, thus 
\begin{equation*}
  0\leq u(z)-\log|g(z)|
\end{equation*}
which implies $u(z)=\log |g(z)|$. It's easy to check that $g\in N$ by definition.
\end{proof}


\begin{theorem}\label{nontangential limit of Nevanlinna}
Let $f\in N$ and $f\not\equiv 0$. Then $f$ has a nontangential limit $f^*(\zeta)$ at almost
every $\zeta\in\mathbb{T}$, and $\log|f^*|\in L^1(\mathbb{T})$.

Moreover, the least harmonic majorant of $\log|f|$ is $P[\mu]$, where
\begin{equation*}
  d\mu=\log|f^*|\,\frac{d\theta}{2\pi}+d\mu_s,
  \qquad \mu_s\perp\frac{d\theta}{2\pi}.
\end{equation*}
\end{theorem}

\subsection{Outer functions and Singular functions}
Before we discuss the main topic of this section, which is the factorization problem of
Nevanlinna functions, we'd like to first introduce two classes of functions: 
\textbf{outer functions} and \textbf{singular functions}.
\begin{definition}\label{def:outer_function}
An analytic function $F$ on $\D$ is an \textbf{outer function} if
\begin{equation*}
  F(z)=C\exp\left(\int K_z(\theta)\,l(\theta)\,\frac{d\theta}{2\pi}\right),
  \qquad z\in\D,
\end{equation*}
for some real-valued $l\in L^1(\mathbb{T})$ and some constant $C\in\mathbb{C}$ with $|C|=1$.
Here $K_z(\theta)=\frac{e^{i\theta}+z}{e^{i\theta}-z}$ is the Herglotz-Riesz kernel.
\end{definition}

\begin{theorem}\label{thm:outer_function_properties}
Let $F$ be an outer function on $\mathbb{D}$. Then:
\begin{itemize}
  \item[(a)] $F$ is analytic on $\D$ and has no zeros in $\D$;
  \item[(b)] $\log|F|=P[l]$ on $\D$;
  \item[(c)] $F\in N$;
  \item[(d)] $F$ has a nontangential limit $F^*(\zeta)$ at almost every $\zeta\in\mathbb{T}$, and
  $|F^*|=e^{l}$ almost everywhere on $\mathbb{T}$;
  \item[(e)] if $F_1,F_2$ are outer functions with $|F_1^*|=|F_2^*|$ almost everywhere on
  $\mathbb{T}$, then $F_1=cF_2$ for some constant $c\in\mathbb{C}$ with $|c|=1$.
\end{itemize}
\end{theorem}
\begin{proof}
Write $F=Ce^{G}$ with
\begin{equation*}
  G(z)=\int K_z(\theta)\,l(\theta)\,\frac{d\theta}{2\pi}.
\end{equation*}

(a) For $|z|\leq r<1$ we have $|K_z(\theta)|\leq\frac{1+r}{1-r}$, and $K_z(\theta)$ is analytic
in $z$. Hence $G$ is analytic on $\D$ and so is
$F=Ce^{G}$. Since the exponential has no zeros, neither does $F$.

(b) As $|C|=1$, $l$ is real-valued and $\text{Re } K_z(\theta)=P_z(\theta)$, note
that $\log |e^G|=\text{Re } G$,

\begin{equation*}
  \log|F(z)|=\text{Re } G(z)=\int P_z(\theta)\,l(\theta)\,\frac{d\theta}{2\pi}=P[l](z).
\end{equation*}

(c) Write $l^+=\max(l,0)$. By Lemma \ref{Properties of Poisson kernel}(a), $P_z>0$, so
$\log|F|=P[l]\leq P[l^+]$, and since $P[l^+]\geq 0$, also $\log^+|F|\leq P[l^+]$. Using
$P_{re^{i\theta}}(\varphi)=P_{re^{i\varphi}}(\theta)$, Fubini's theorem and Lemma
\ref{Properties of Poisson kernel}(a) again,
\begin{equation*}
  \frac{1}{2\pi}\int_0^{2\pi}\log^+|F(re^{i\theta})|\,d\theta
  \leq\frac{1}{2\pi}\int_0^{2\pi}P[l^+](re^{i\theta})\,d\theta
  =\frac{1}{2\pi}\int_0^{2\pi}l^+(\varphi)\,d\varphi
  \leq\|l\|_{L^1}
\end{equation*}
for every $0<r<1$. Hence $F\in N$.

(d) By (c) and $F\not\equiv0$, Theorem \ref{nontangential limit of Nevanlinna} shows that $F$ has
a nontangential limit $F^*$ almost everywhere. By (b) and Corollary
\ref{cor:poisson_nontangential_ae}, $\log|F|=P[l]$ has nontangential limit $l$ almost
everywhere. At almost every $\zeta\in\mathbb{T}$ both limits exist, and then
\begin{equation*}
  |F^*(\zeta)|=\lim_{\substack{z\to\zeta\\ z\in\Gamma_\alpha(\zeta)}}e^{\log|F(z)|}=e^{l(\zeta)}.
\end{equation*}

(e) Let $F_j=C_je^{G_j}$, where $G_j$ is defined from $l_j$ as above, $j=1,2$. By (d),
$e^{l_1}=|F_1^*|=|F_2^*|=e^{l_2}$ almost everywhere, so $l_1=l_2$ almost everywhere. Hence
$G_1=G_2$ on $\D$, and $F_1=cF_2$ with $c=C_1/C_2$, $|c|=1$.
\end{proof}

\begin{remark}
One can construct an outer function as follow: given $h:\mathbb{T}\rightarrow [0,\infty)$
s.t. $\log h\in L^1(\mathbb{T})$
(think of it as the modulos of the boundary function of Nevanlinna functions over $\mathbb{D}$, 
see \ref{nontangential limit of Nevanlinna}), then set $u(z)=P[\log h](z)$ and $v(z)$ to be one of the
harmonic conjugates of $u(z)$. Then the function $F(z):= e^{u(z)+i v(z)}$ is an outer function,
and it is determined by $h(\theta)$ except for a unimodular constant factor due to the choice
of the conjugate $v(z)$. 
\end{remark}

\begin{definition}
An analytic function $S$ on $\D$ is a \textbf{singular function} if
\begin{equation*}
  S(z)=\exp\left(-\int K_z(\theta)\,d\mu(\theta)\right),
  \qquad z\in\D,
\end{equation*}
for some finite singular positive measure $\mu$ on $\mathbb{T}$.
\end{definition}

\begin{lemma}\label{lem:log_singular_function}
For singular function $S(z)=\exp\left(-\int K_z(\theta)\,d\mu(\theta)\right)$
 on $\mathbb{D}$, we have $\log|S|=-P[\mu]$ on $\mathbb{D}$.
\end{lemma}
\begin{proof}
Given $z\in\D$, as $|K_z(\theta)|\leq\frac{1+|z|}{1-|z|}$ and $\mu$ is finite, 
$\int K_z(\theta)\,d\mu(\theta)$ converges. As $\log|e^w|=\text{Re } w$, $\mu$ is a positive
measure and $\text{Re } K_z(\theta)=P_z(\theta)$,
\begin{equation*}
  \log|S(z)|=-\text{Re }\int K_z(\theta)\,d\mu(\theta)
  =-\int P_z(\theta)\,d\mu(\theta)=-P[\mu](z).
\end{equation*}
\end{proof}

\begin{lemma}\label{lem:singular_function_in_N}
If $S$ is a singular function on $\mathbb{D}$, then $S\in N$.
\end{lemma}
\begin{proof}
By Lemma \ref{lem:log_singular_function} and Lemma \ref{Properties of Poisson kernel}(a),
$\log|S|=-P[\mu]\leq 0$ on $\D$, as $\mu\geq 0$ and $P_z>0$. Hence $\log^+|S|=0$ on $\D$,
so $S\in N$.
\end{proof}

\begin{theorem}\label{equivalent characterization of singular function}
Let $S$ be an analytic function on $\D$. Then $S$ is a singular function if and only if $S$ has the
following properties:
\begin{itemize}
  \item[(a)] $S$ has no zeros in $\D$;
  \item[(b)] $|S(z)|\leq 1$ for all $z\in\D$;
  \item[(c)] $S$ has a nontangential limit $S^*(\zeta)$ at almost every $\zeta\in\mathbb{T}$, and
  $|S^*|=1$ almost everywhere on $\mathbb{T}$;
  \item[(d)] $S(0)>0$.
\end{itemize}
\end{theorem}

\begin{proof}[Proof of one direction]
($\Rightarrow$) Let $S(z)=\exp\left(-\int K_z(\theta)\,d\mu(\theta)\right)$ with $\mu\geq 0$
finite and singular.
(a) holds as exponential has no zeros.
(b) was shown in the proof of Lemma \ref{lem:singular_function_in_N}.
(c) By Lemma \ref{lem:singular_function_in_N} and Theorem \ref{nontangential limit of Nevanlinna},
$S$ has a nontangential limit $S^*$ almost everywhere. By Lemma \ref{lem:log_singular_function}
and Theorem \ref{thm:singular_poisson_nontangential_zero}, $\log|S|=-P[\mu]$ has nontangential
limit $0$ almost everywhere. So at a.e. $\zeta\in\mathbb{T}$,
\begin{equation*}
  |S^*(\zeta)|=\lim_{\substack{z\to\zeta\\ z\in\Gamma_\alpha(\zeta)}}e^{-P[\mu](z)}=1.
\end{equation*}
(d) holds because $K_0(\theta)=1$ gives $S(0)=e^{-\mu(\mathbb{T})}>0$.

We will skip the proof of ($\Leftarrow$) as it's not needed in the project.
\end{proof}



\subsection{Canonical factorization theorem of Nevanlinna funcitons}
Now we'd like to discuss the \textit{factorization problem} of Nevanlinna functions.
Given $f\in N$ with $f\not\equiv 0$, consider $B$: Blaschke product formed by
its zeros and $g=f/B$. By Lemma \ref{least harmonic majorant of log|f| for Nevanlinna f},
we know that $\log|g|$ is the least harmonic majorant of $\log |f|$. Recall 
Theorem \ref{characterization of Nevanlinna using log |f|}, we can write 
$\log |g|=P[\mu]$ for some finite measure $\mu$. By Theorem \ref{nontangential limit of Nevanlinna},
we can write:
\begin{equation*}
  d\mu=\log|f^*|\,\frac{d\theta}{2\pi}+d\mu_s,
  \qquad \mu_s\perp\frac{d\theta}{2\pi}.
\end{equation*}
Now take Jordan decomposition on singular measure
 $\mu_s$, we get the unique decomposition:
 \begin{equation*}
  \mu_s=-\mu_1+\mu_2
 \end{equation*}
 where $\mu_1$ and $\mu_2$ are positive singular measure that are mutually singular to each other. Thus we 
write:
\begin{equation*}
  d\mu=\log|f^*|\,\frac{d\theta}{2\pi}-d\mu_1+d\mu_2.
\end{equation*}

Now consider the \textbf{outer function}
\begin{equation*}
  \tilde{F}(z)=\exp\left(\int K_z(\theta)\log |f^*(e^{i\theta})|\frac{d\theta}{2\pi}\right)
\end{equation*}
Here $K_z(\theta)=\frac{e^{i\theta}+z}{e^{i\theta}-z}$ is the Herglotz-Riesz kernel. Moreover, consider
the \textbf{singular functions}:
\begin{equation*}
  S_j(z)=\exp\left(-\int K_z(\theta)d\mu_j(\theta)\right),\quad j=1,2.
\end{equation*}
Recall that $\text{Re } K_z(\theta)=P_z(\theta)$, which is the Poisson kernel, and
$\log |e^w|=\text{Re } w$, so 
\begin{equation*}
  \log |\tilde{F}(z)|=P[\log |f|](z),\; \log |S_j(z)|=- P[\mu_j](z),\quad j=1,2.
\end{equation*}
Therefore, 
\begin{align*}
  \log |g(z)|=P[\mu](z)&=P[\log |f|](z)+P[\mu_1](z)-P[\mu_2](z)\\
  &= \log |\tilde{F}(z)|+\log |S_1(z)|-\log|S_2(z)|.
\end{align*}
Then we have:
\begin{equation*}
  \log \left|\frac{g(z) S_2(z)}{\tilde{F}(z)S_1(z)}\right|=0
\end{equation*}
which takes $\exp$ on both sides yields:
\begin{equation*}
  \left|\frac{f(z)S_2(z)}{B(z)\tilde{F}(z)S_1(z)}\right|=1,\quad \forall z\in \mathbb{D}.
\end{equation*}
By maximum modulus principle, 
\begin{equation*}
  \frac{f(z)S_2(z)}{B(z)\tilde{F}(z)S_1(z)}=C,\quad C\in\mathbb{C},\; |C|=1.
\end{equation*}
Set $F=C\tilde{F}$. By Definition \ref{def:outer_function}, $F$ is again an outer function, and
$f=BFS_1/S_2$.
This gives most of the \textbf{canonical factorization theorem}:
\begin{theorem}
Let $f\in N$, $f\not\equiv 0$. Then
\begin{equation*}
  f(z)=B(z)F(z)S_1(z)/S_2(z),
\end{equation*}
where $B$ is Blaschke product formed by zeros of $f$,
$F$ is an outer function, and $S_1,S_2$ are
singular functions whose positive singular measures $\mu_1,\mu_2$ are mutually singular.
This factorization is unique.
Conversely, every function of this form is in $N$.
\end{theorem}
\begin{proof}
The existence part is given as above.

For uniqueness, the Blaschke product $B$ is determined by the zeros of $f$. Suppose
$f=BFS_1/S_2=BF'S_1'/S_2'$, where $S_1,S_2$ come from $\mu_1,\mu_2$ and $S_1',S_2'$ from
$\nu_1,\nu_2$. Since the zeros of $B\not\equiv 0$ are isolated, $FS_1/S_2=F'S_1'/S_2'$ on $\D$
minus a discrete set, hence on $\D$ by continuity. By Theorem
\ref{equivalent characterization of singular function}(c), $|S_1^*|=|S_2^*|=|S_1'^*|=|S_2'^*|=1$
almost everywhere, so $|F^*|=|F'^*|$ almost everywhere, and by Theorem
\ref{thm:outer_function_properties}(e), $F=cF'$ for a constant $c$ with $|c|=1$. Thus
$cS_1/S_2=S_1'/S_2'$, and taking $\log|\cdot|$, Lemma \ref{lem:log_singular_function} gives
\begin{equation*}
  P[\mu_2-\mu_1]=P[\nu_2-\nu_1]\quad\text{on }\D.
\end{equation*}
Since a finite measure on $\mathbb{T}$ is determined by its Poisson integral,
$\mu_2-\mu_1=\nu_2-\nu_1$. Both $(\mu_2,\mu_1)$ and $(\nu_2,\nu_1)$ are pairs of mutually
singular positive measures, so by uniqueness of the Jordan decomposition, $\mu_1=\nu_1$ and
$\mu_2=\nu_2$. Hence $S_j=S_j'$ for $j=1,2$, so $c=1$ and $F=F'$.

For converse part Since $N$ is closed under multiplication by Remark \ref{rem:N_algebra}, it suffices to show that
$B,F,S_1,1/S_2\in N$. We have $B\in H^\infty\subset N$ by Theorem \ref{Blaschke product},
$F\in N$ by Theorem \ref{thm:outer_function_properties}(c), and $S_1\in N$ by Lemma
\ref{lem:singular_function_in_N}. Likewise $S_2\in N$, and $S_2$ has no zeros in $\D$, so
$1/S_2\in N$ by Remark \ref{rem:N_algebra}.
\end{proof}
.

\chapter{Hardy spaces and Nevanlinna class}
\label{sec:hardy_nevanlinna}
This chapter gives an introduction on the mathematical backgrounds of
the main objects of this thesis: Hardy spaces and 
Nevanlinna class, and discuss the ultimate goal to formalize in
this project: canonical factorization theorem of Nevanlinna functions. This chapter is
mainly based on \cite{garnett2006bounded}, II.1 - II.5.

\section{Hardy spaces}
One of the main objects of this thesis is the \textbf{Hardy space}, which is a subspace of the space 
$\mathcal{O}(U)$ of analytic functions
on a `certain domain' $U\subset \mathbb{C}$ which are defined by 
`certain growth conditions'. Roughly speaking,
`certain growth conditions' means that the boundary growth is controlled in 
an averaged $L^p$ sense. We usually consider
$U$ to be the unit disc $\D:=\{z\in\C: |z|<1\}$ or the upper half plane 
$\HH:=\{z\in\C: \im z>0\}$.

\begin{definition}
Let $0<p<\infty$. We define the \textbf{Hardy space over the unit disc $\D$} as:
\begin{equation*}
  H^p(\D):=\left\{ f\in \mathcal{O}(\D):\| f\|_{H^p}^p := \sup_{0<r<1}\frac{1}{2\pi}\int |f(re^{i\theta})|^p d
  \theta <\infty\right\}
\end{equation*}
If $p=\infty$, we define $H^\infty(\D)$ to be the space of bounded analytic functions on $\D$, i.e.
\begin{equation*}
  H^\infty(\D):=\{f\in \mathcal{O}(\D): \|f\|_{H^\infty}:=\sup_{z\in \D}|f(z)|<\infty\}
\end{equation*}
\end{definition}

\begin{definition}
Let $0<p<\infty$. We define the \textbf{Hardy space over the upper half plane $\HH$} as:
\begin{equation*}
  H^p(dt):=\left\{f\in \mathcal{O}(\HH): \|f\|_{H^p}^p:=\sup_{y>0}\int |f(x+iy)|^p dx <\infty\right\}
\end{equation*}
If $p=\infty$, we define $H^\infty(dt)$ to be the space of bounded analytic functions on $\HH$, i.e.
\begin{equation*}
  H^\infty(dt):=\{f\in \mathcal{O}(\HH): \|f\|_{H^\infty}:=\sup_{z\in \HH}|f(z)|<\infty\}
\end{equation*}
\end{definition}
These two definitions are closely related and many results holds true on $H^p(\D)$ will be also holds
true on $H^p(dt)$. In this thesis, we only focus on the formalization for $H^p(\D)$.

\begin{remark}
One can extend the definition of Hardy space over general (bounded) domain in $\mathbb{C}$ as follow:
\begin{definition}
Let $0<p<\infty$ and $U\subset \mathbb{C}$ be a (bounded) domain. We say an analytic function $f\in \mathcal{O}(U)$
over $U$ is a $H^p$ function if the function $|f|^p$ has a harmonic majorant, and we define $\|f\|^p_{H^p}$
to be the value of the least harmonic majorant at $z=0$. For $p=\infty$, $f\in \mathcal{O}(U)$ is
$H^\infty$ if it is a bounded analytic function over $U$ and $\|f\|_{H^\infty}:=\sup_{z\in U}|f(z)|$.
\end{definition}
Theorem \ref{Harmonic majorant criterion} shows that this is an equivalent definition with above definitions over $\D$.
\end{remark}


The Hardy space $H^p(\D)$ is a normed space for $p\geq 1$ with the
norm $\|\cdot\|_{H^p}$. Moreover, it is a metric space for all $0<p\leq \infty$,
with the metric $d(f,g)=\|f-g\|_{H^p}^p$ for $0<p<1$ and the metric induced by
the norm for $p\geq 1$. Moreover, we have:
\begin{theorem}
For $0<p\leq \infty$, $H^p$ is a complete metric space.
\end{theorem}
Similar as $L^p$ space, there is monotonicity for $H^p$ space as well:
\begin{theorem}
For $0<p\leq q\leq \infty$, we have $H^q\leq H^p$ as subspace.
\end{theorem}


\section{Blaschke product}
\begin{definition}
Given a sequence $\{z_n\}$ on $\mathbb{D}$, we say that this sequence satisfies
\textbf{Blaschke's condition} if $\sum (1-|z_n|)<\infty$. Given an analytic function $f$ on 
$\mathbb{D}$, we call this function $f$ satisfies \textbf{Blaschke's condition} if its
zeros $\{z_n\}$ satisfies Blaschke's condition.
\end{definition}

\begin{theorem}\label{Blaschke condition from harmonic majorant}
Let $f\not\equiv 0$ be an analytic function on $\mathbb{D}$ with zeros $\{z_n\}$. If
$\log|f|$ has a harmonic majorant, then $f$ satisfies Blaschke's condition:
\begin{equation*}
  \sum(1-|z_n|)<\infty.
\end{equation*}
If moreover $f(0)\neq 0$ and $u$ is the least harmonic majorant of $\log|f|$, then
\begin{equation*}
  \sum(1-|z_n|)\leq u(0)-\log|f(0)|.
\end{equation*}
\end{theorem}

\begin{theorem}\label{Blaschke product}
Let $\{z_n\}$ be a sequence in $\mathbb{D}$ satisfying Blaschke's condition, and let $m$ be the
number of $z_n$ equal to $0$. Then the \textbf{Blaschke product}
\begin{equation}\label{eq:Blaschke product}
  B(z)=z^m\prod_{|z_n|\neq 0}\left(-\frac{\overline{z_n}}{|z_n|}\right)
  \frac{z-z_n}{1-\overline{z_n}z}
\end{equation}
converges on $\mathbb{D}$. The function $B$ is in $H^\infty(\mathbb{D})$ with $|B(z)|\leq 1$,
and the zeros of $B$ are precisely the points $z_n$, each zero having multiplicity equal to
the number of times it occurs in $\{z_n\}$. Moreover,
\begin{equation*}
  |B(e^{i\theta})|=1
\end{equation*}
for almost every $\theta$, where $B(e^{i\theta})=\lim_{r\rightarrow 1}B(re^{i\theta})$ is the
radial limit of $B$.
\end{theorem}

Here's a characterization of Blaschke products in terms of harmonic majorants:
\begin{theorem}\label{characterization of Blaschke products}
Let $f\in H^\infty(\mathbb{D})$ with $\|f\|_{H^\infty}\leq 1$. The followings are equivalent:
\begin{itemize}
  \item[(a)] $f(z)=\lambda B(z)$ where $\lambda$ is constant with $|\lambda|=1$ and $B$ is a 
  Blaschke product;
  \item[(b)] $\lim_{r\rightarrow 1}\int \log |f(re^{i\theta})|d\theta/(2\pi)=0$;
  \item[(c)] The least harmonic majorant of $\log |f|$ is 0.
\end{itemize}
\end{theorem}

\section{Boundary functions of Hardy functions}
Every $H^p$ function has a boundary function in $L^p(\mathbb{T})$, and for $p\geq 1$ it can be
recovered from its boundary function via Poisson integral.
\begin{theorem}\label{thm:Hp_boundary_function}
Let $0<p\leq\infty$ and $f\in H^p(\D)$. Then $f$ has a nontangential limit $f^*(\zeta)$ at
almost every $\zeta\in\mathbb{T}$, and $f^*\in L^p(\mathbb{T}, \frac{d\theta}{2\pi})$ with
\begin{equation*}
  \|f^*\|_{L^p}=\|f\|_{H^p}.
\end{equation*}
Moreover:
\begin{itemize}
  \item[(a)] if $0<p<\infty$, then
  \begin{equation*}
    \lim_{r\rightarrow 1}\|f(re^{i\cdot})-f^*\|_{L^p}=0;
  \end{equation*}
  \item[(b)] if $1\leq p\leq\infty$, then $f=P[f^*]$.
\end{itemize}
\end{theorem}




\section{Nevanlinna class} 
\begin{definition}
An analytic function $f(z)$ on $\D$ is in the \textbf{Nevanlinna class} $N$ if 
\begin{equation*}
  \sup_r \int \log^+|f(r e^{i\theta})|\frac{d\theta}{2\pi}<\infty
\end{equation*}
\end{definition}
As $\log^+|f(z)|\leq (1/p)|f(z)|^p$, so $H^p\subset N$ for all $0<p\leq \infty$, i.e. Nevanlinna
class encompasses all Hardy spaces $H^p$. 

By Theorem \ref{Harmonic majorant criterion}, we can formulate the definition
in an equivalent definition:
\begin{definition}
An analytic function $f(z)$ on $\D$ is in the \textbf{Nevanlinna class} $N$ if 
the subharmonic function $\log^+|f(z)|$ has a harmonic majorant.
\end{definition}

\begin{remark}\label{rem:N_algebra}
Compared to $H^p(\D) \subset \mathcal{O}(\D)$,
which is just a subspace of $\mathcal{O}(\D)$ for $p<\infty$, $N\subset \mathcal{O}(\D)$ is
also a subalgebra of $\mathcal{O}(\D)$: if $f,g\in N$, then $fg\in N$ as
\begin{equation*}
  \log^+|fg|\leq \log^+|f|+\log^+|g|.
\end{equation*}
Moreover, if $f\in N$ has no zeros in $\D$, then $1/f\in N$. Indeed, let $U$ be a harmonic
majorant of $\log^+|f|$. Since $f$ has no zeros, $\log|f|$ is harmonic on $\D$, and
\begin{equation*}
  \log^+|1/f|=\log^+|f|-\log|f|\leq U-\log|f|,
\end{equation*}
so $\log^+|1/f|$ has the harmonic majorant $U-\log|f|$.
The closure of multiplication and inversion will sort of provide a more natural place for us to 
describe `factorization problem'.
\end{remark}

\begin{theorem}\label{characterization of Nevanlinna using log |f|}
Let $f(z)$ be an analytic function on $\mathbb{D}$ s.t. $f\not\equiv 0$. 
Then $f\in N$ if and only if $\log|f(z)|$ has least harmonic majorant
the Poisson integral of some finite measure on $\mathbb{T}$.
\end{theorem}
\begin{proof}

\end{proof}


\begin{lemma}\label{least harmonic majorant of log|f| for Nevanlinna f}
Let $f\in N$ and $f\not\equiv 0$. Let $B$ be the Blaschke product formed from zeros of $f$. Then:
\begin{itemize}
  \item $B$ converges;
  \item $g=f/B\in N$;
  \item $\log |g|$ is the least harmonic majorant of $\log |f|$.
\end{itemize}
\end{lemma}
\begin{proof}[Sketch of proof]
As $f\in N$, $\log |f|$ has a harmonic majorant, thus by Theorem \ref{Blaschke condition from harmonic majorant},
$f$ satisfies Blaschke's condition, thus $B$ converges. Now let $u$ be the least harmonic majorant
of $\log |f|$. Recall in Example \ref{log |f| is subharmonic}, $g$ as a zero-free analytic function
over $\mathbb{D}$, $\log |g|$ is harmonic on $\mathbb{D}$. As $|B(z)|\leq 1$, we have 
$\log |f(z)|\leq \log |g(z)|$, thus 
\begin{equation*}
  u(z)\leq \log|g(z)|.
\end{equation*}
On the other hand, 
\begin{equation*}
  \log |B(z)|=\log |f(z)|-\log|g(z)|\leq u(z)-\log |g(z)|
\end{equation*}
Recall Theorem \ref{characterization of Blaschke products}, the least harmonic majorant of $\log|B|$
is 0, thus 
\begin{equation*}
  0\leq u(z)-\log|g(z)|
\end{equation*}
which implies $u(z)=\log |g(z)|$. It's easy to check that $g\in N$ by definition.
\end{proof}


\begin{theorem}\label{nontangential limit of Nevanlinna}
Let $f\in N$ and $f\not\equiv 0$. Then $f$ has a nontangential limit $f^*(\zeta)$ at almost
every $\zeta\in\mathbb{T}$, and $\log|f^*|\in L^1(\mathbb{T})$.

Moreover, the least harmonic majorant of $\log|f|$ is $P[\mu]$, where
\begin{equation*}
  d\mu=\log|f^*|\,\frac{d\theta}{2\pi}+d\mu_s,
  \qquad \mu_s\perp\frac{d\theta}{2\pi}.
\end{equation*}
\end{theorem}

\subsection{Outer functions and Singular functions}
Before we discuss the main topic of this section, which is the factorization problem of
Nevanlinna functions, we'd like to first introduce two classes of functions: 
\textbf{outer functions} and \textbf{singular functions}.
\begin{definition}\label{def:outer_function}
An analytic function $F$ on $\D$ is an \textbf{outer function} if
\begin{equation*}
  F(z)=C\exp\left(\int K_z(\theta)\,l(\theta)\,\frac{d\theta}{2\pi}\right),
  \qquad z\in\D,
\end{equation*}
for some real-valued $l\in L^1(\mathbb{T})$ and some constant $C\in\mathbb{C}$ with $|C|=1$.
Here $K_z(\theta)=\frac{e^{i\theta}+z}{e^{i\theta}-z}$ is the Herglotz-Riesz kernel.
\end{definition}

\begin{theorem}\label{thm:outer_function_properties}
Let $F$ be an outer function on $\mathbb{D}$. Then:
\begin{itemize}
  \item[(a)] $F$ is analytic on $\D$ and has no zeros in $\D$;
  \item[(b)] $\log|F|=P[l]$ on $\D$;
  \item[(c)] $F\in N$;
  \item[(d)] $F$ has a nontangential limit $F^*(\zeta)$ at almost every $\zeta\in\mathbb{T}$, and
  $|F^*|=e^{l}$ almost everywhere on $\mathbb{T}$;
  \item[(e)] if $F_1,F_2$ are outer functions with $|F_1^*|=|F_2^*|$ almost everywhere on
  $\mathbb{T}$, then $F_1=cF_2$ for some constant $c\in\mathbb{C}$ with $|c|=1$.
\end{itemize}
\end{theorem}
\begin{proof}
Write $F=Ce^{G}$ with
\begin{equation*}
  G(z)=\int K_z(\theta)\,l(\theta)\,\frac{d\theta}{2\pi}.
\end{equation*}

(a) For $|z|\leq r<1$ we have $|K_z(\theta)|\leq\frac{1+r}{1-r}$, and $K_z(\theta)$ is analytic
in $z$. Hence $G$ is analytic on $\D$ and so is
$F=Ce^{G}$. Since the exponential has no zeros, neither does $F$.

(b) As $|C|=1$, $l$ is real-valued and $\text{Re } K_z(\theta)=P_z(\theta)$, note
that $\log |e^G|=\text{Re } G$,

\begin{equation*}
  \log|F(z)|=\text{Re } G(z)=\int P_z(\theta)\,l(\theta)\,\frac{d\theta}{2\pi}=P[l](z).
\end{equation*}

(c) Write $l^+=\max(l,0)$. By Lemma \ref{Properties of Poisson kernel}(a), $P_z>0$, so
$\log|F|=P[l]\leq P[l^+]$, and since $P[l^+]\geq 0$, also $\log^+|F|\leq P[l^+]$. Using
$P_{re^{i\theta}}(\varphi)=P_{re^{i\varphi}}(\theta)$, Fubini's theorem and Lemma
\ref{Properties of Poisson kernel}(a) again,
\begin{equation*}
  \frac{1}{2\pi}\int_0^{2\pi}\log^+|F(re^{i\theta})|\,d\theta
  \leq\frac{1}{2\pi}\int_0^{2\pi}P[l^+](re^{i\theta})\,d\theta
  =\frac{1}{2\pi}\int_0^{2\pi}l^+(\varphi)\,d\varphi
  \leq\|l\|_{L^1}
\end{equation*}
for every $0<r<1$. Hence $F\in N$.

(d) By (c) and $F\not\equiv0$, Theorem \ref{nontangential limit of Nevanlinna} shows that $F$ has
a nontangential limit $F^*$ almost everywhere. By (b) and Corollary
\ref{cor:poisson_nontangential_ae}, $\log|F|=P[l]$ has nontangential limit $l$ almost
everywhere. At almost every $\zeta\in\mathbb{T}$ both limits exist, and then
\begin{equation*}
  |F^*(\zeta)|=\lim_{\substack{z\to\zeta\\ z\in\Gamma_\alpha(\zeta)}}e^{\log|F(z)|}=e^{l(\zeta)}.
\end{equation*}

(e) Let $F_j=C_je^{G_j}$, where $G_j$ is defined from $l_j$ as above, $j=1,2$. By (d),
$e^{l_1}=|F_1^*|=|F_2^*|=e^{l_2}$ almost everywhere, so $l_1=l_2$ almost everywhere. Hence
$G_1=G_2$ on $\D$, and $F_1=cF_2$ with $c=C_1/C_2$, $|c|=1$.
\end{proof}

\begin{remark}
One can construct an outer function as follow: given $h:\mathbb{T}\rightarrow [0,\infty)$
s.t. $\log h\in L^1(\mathbb{T})$
(think of it as the modulos of the boundary function of Nevanlinna functions over $\mathbb{D}$, 
see \ref{nontangential limit of Nevanlinna}), then set $u(z)=P[\log h](z)$ and $v(z)$ to be one of the
harmonic conjugates of $u(z)$. Then the function $F(z):= e^{u(z)+i v(z)}$ is an outer function,
and it is determined by $h(\theta)$ except for a unimodular constant factor due to the choice
of the conjugate $v(z)$. 
\end{remark}

\begin{definition}
An analytic function $S$ on $\D$ is a \textbf{singular function} if
\begin{equation*}
  S(z)=\exp\left(-\int K_z(\theta)\,d\mu(\theta)\right),
  \qquad z\in\D,
\end{equation*}
for some finite singular positive measure $\mu$ on $\mathbb{T}$.
\end{definition}

\begin{lemma}\label{lem:log_singular_function}
For singular function $S(z)=\exp\left(-\int K_z(\theta)\,d\mu(\theta)\right)$
 on $\mathbb{D}$, we have $\log|S|=-P[\mu]$ on $\mathbb{D}$.
\end{lemma}
\begin{proof}
Given $z\in\D$, as $|K_z(\theta)|\leq\frac{1+|z|}{1-|z|}$ and $\mu$ is finite, 
$\int K_z(\theta)\,d\mu(\theta)$ converges. As $\log|e^w|=\text{Re } w$, $\mu$ is a positive
measure and $\text{Re } K_z(\theta)=P_z(\theta)$,
\begin{equation*}
  \log|S(z)|=-\text{Re }\int K_z(\theta)\,d\mu(\theta)
  =-\int P_z(\theta)\,d\mu(\theta)=-P[\mu](z).
\end{equation*}
\end{proof}

\begin{lemma}\label{lem:singular_function_in_N}
If $S$ is a singular function on $\mathbb{D}$, then $S\in N$.
\end{lemma}
\begin{proof}
By Lemma \ref{lem:log_singular_function} and Lemma \ref{Properties of Poisson kernel}(a),
$\log|S|=-P[\mu]\leq 0$ on $\D$, as $\mu\geq 0$ and $P_z>0$. Hence $\log^+|S|=0$ on $\D$,
so $S\in N$.
\end{proof}

\begin{theorem}\label{equivalent characterization of singular function}
Let $S$ be an analytic function on $\D$. Then $S$ is a singular function if and only if $S$ has the
following properties:
\begin{itemize}
  \item[(a)] $S$ has no zeros in $\D$;
  \item[(b)] $|S(z)|\leq 1$ for all $z\in\D$;
  \item[(c)] $S$ has a nontangential limit $S^*(\zeta)$ at almost every $\zeta\in\mathbb{T}$, and
  $|S^*|=1$ almost everywhere on $\mathbb{T}$;
  \item[(d)] $S(0)>0$.
\end{itemize}
\end{theorem}

\begin{proof}[Proof of one direction]
($\Rightarrow$) Let $S(z)=\exp\left(-\int K_z(\theta)\,d\mu(\theta)\right)$ with $\mu\geq 0$
finite and singular.
(a) holds as exponential has no zeros.
(b) was shown in the proof of Lemma \ref{lem:singular_function_in_N}.
(c) By Lemma \ref{lem:singular_function_in_N} and Theorem \ref{nontangential limit of Nevanlinna},
$S$ has a nontangential limit $S^*$ almost everywhere. By Lemma \ref{lem:log_singular_function}
and Theorem \ref{thm:singular_poisson_nontangential_zero}, $\log|S|=-P[\mu]$ has nontangential
limit $0$ almost everywhere. So at a.e. $\zeta\in\mathbb{T}$,
\begin{equation*}
  |S^*(\zeta)|=\lim_{\substack{z\to\zeta\\ z\in\Gamma_\alpha(\zeta)}}e^{-P[\mu](z)}=1.
\end{equation*}
(d) holds because $K_0(\theta)=1$ gives $S(0)=e^{-\mu(\mathbb{T})}>0$.

We will skip the proof of ($\Leftarrow$) as it's not needed in the project.
\end{proof}



\subsection{Canonical factorization theorem of Nevanlinna funcitons}
Now we'd like to discuss the \textit{factorization problem} of Nevanlinna functions.
Given $f\in N$ with $f\not\equiv 0$, consider $B$: Blaschke product formed by
its zeros and $g=f/B$. By Lemma \ref{least harmonic majorant of log|f| for Nevanlinna f},
we know that $\log|g|$ is the least harmonic majorant of $\log |f|$. Recall 
Theorem \ref{characterization of Nevanlinna using log |f|}, we can write 
$\log |g|=P[\mu]$ for some finite measure $\mu$. By Theorem \ref{nontangential limit of Nevanlinna},
we can write:
\begin{equation*}
  d\mu=\log|f^*|\,\frac{d\theta}{2\pi}+d\mu_s,
  \qquad \mu_s\perp\frac{d\theta}{2\pi}.
\end{equation*}
Now take Jordan decomposition on singular measure
 $\mu_s$, we get the unique decomposition:
 \begin{equation*}
  \mu_s=-\mu_1+\mu_2
 \end{equation*}
 where $\mu_1$ and $\mu_2$ are positive singular measure that are mutually singular to each other. Thus we 
write:
\begin{equation*}
  d\mu=\log|f^*|\,\frac{d\theta}{2\pi}-d\mu_1+d\mu_2.
\end{equation*}

Now consider the \textbf{outer function}
\begin{equation*}
  \tilde{F}(z)=\exp\left(\int K_z(\theta)\log |f^*(e^{i\theta})|\frac{d\theta}{2\pi}\right)
\end{equation*}
Here $K_z(\theta)=\frac{e^{i\theta}+z}{e^{i\theta}-z}$ is the Herglotz-Riesz kernel. Moreover, consider
the \textbf{singular functions}:
\begin{equation*}
  S_j(z)=\exp\left(-\int K_z(\theta)d\mu_j(\theta)\right),\quad j=1,2.
\end{equation*}
Recall that $\text{Re } K_z(\theta)=P_z(\theta)$, which is the Poisson kernel, and
$\log |e^w|=\text{Re } w$, so 
\begin{equation*}
  \log |\tilde{F}(z)|=P[\log |f|](z),\; \log |S_j(z)|=- P[\mu_j](z),\quad j=1,2.
\end{equation*}
Therefore, 
\begin{align*}
  \log |g(z)|=P[\mu](z)&=P[\log |f|](z)+P[\mu_1](z)-P[\mu_2](z)\\
  &= \log |\tilde{F}(z)|+\log |S_1(z)|-\log|S_2(z)|.
\end{align*}
Then we have:
\begin{equation*}
  \log \left|\frac{g(z) S_2(z)}{\tilde{F}(z)S_1(z)}\right|=0
\end{equation*}
which takes $\exp$ on both sides yields:
\begin{equation*}
  \left|\frac{f(z)S_2(z)}{B(z)\tilde{F}(z)S_1(z)}\right|=1,\quad \forall z\in \mathbb{D}.
\end{equation*}
By maximum modulus principle, 
\begin{equation*}
  \frac{f(z)S_2(z)}{B(z)\tilde{F}(z)S_1(z)}=C,\quad C\in\mathbb{C},\; |C|=1.
\end{equation*}
Set $F=C\tilde{F}$. By Definition \ref{def:outer_function}, $F$ is again an outer function, and
$f=BFS_1/S_2$.
This gives most of the \textbf{canonical factorization theorem}:
\begin{theorem}
Let $f\in N$, $f\not\equiv 0$. Then
\begin{equation*}
  f(z)=B(z)F(z)S_1(z)/S_2(z),
\end{equation*}
where $B$ is Blaschke product formed by zeros of $f$,
$F$ is an outer function, and $S_1,S_2$ are
singular functions whose positive singular measures $\mu_1,\mu_2$ are mutually singular.
This factorization is unique.
Conversely, every function of this form is in $N$.
\end{theorem}
\begin{proof}
The existence part is given as above.

For uniqueness, the Blaschke product $B$ is determined by the zeros of $f$. Suppose
$f=BFS_1/S_2=BF'S_1'/S_2'$, where $S_1,S_2$ come from $\mu_1,\mu_2$ and $S_1',S_2'$ from
$\nu_1,\nu_2$. Since the zeros of $B\not\equiv 0$ are isolated, $FS_1/S_2=F'S_1'/S_2'$ on $\D$
minus a discrete set, hence on $\D$ by continuity. By Theorem
\ref{equivalent characterization of singular function}(c), $|S_1^*|=|S_2^*|=|S_1'^*|=|S_2'^*|=1$
almost everywhere, so $|F^*|=|F'^*|$ almost everywhere, and by Theorem
\ref{thm:outer_function_properties}(e), $F=cF'$ for a constant $c$ with $|c|=1$. Thus
$cS_1/S_2=S_1'/S_2'$, and taking $\log|\cdot|$, Lemma \ref{lem:log_singular_function} gives
\begin{equation*}
  P[\mu_2-\mu_1]=P[\nu_2-\nu_1]\quad\text{on }\D.
\end{equation*}
Since a finite measure on $\mathbb{T}$ is determined by its Poisson integral,
$\mu_2-\mu_1=\nu_2-\nu_1$. Both $(\mu_2,\mu_1)$ and $(\nu_2,\nu_1)$ are pairs of mutually
singular positive measures, so by uniqueness of the Jordan decomposition, $\mu_1=\nu_1$ and
$\mu_2=\nu_2$. Hence $S_j=S_j'$ for $j=1,2$, so $c=1$ and $F=F'$.

For converse part Since $N$ is closed under multiplication by Remark \ref{rem:N_algebra}, it suffices to show that
$B,F,S_1,1/S_2\in N$. We have $B\in H^\infty\subset N$ by Theorem \ref{Blaschke product},
$F\in N$ by Theorem \ref{thm:outer_function_properties}(c), and $S_1\in N$ by Lemma
\ref{lem:singular_function_in_N}. Likewise $S_2\in N$, and $S_2$ has no zeros in $\D$, so
$1/S_2\in N$ by Remark \ref{rem:N_algebra}.
\end{proof}
.

\chapter{Hardy spaces and Nevanlinna class}
\label{sec:hardy_nevanlinna}
This chapter gives an introduction on the mathematical backgrounds of
the main objects of this thesis: Hardy spaces and 
Nevanlinna class, and discuss the ultimate goal to formalize in
this project: canonical factorization theorem of Nevanlinna functions. This chapter is
mainly based on \cite{garnett2006bounded}, II.1 - II.5.

\section{Hardy spaces}
One of the main objects of this thesis is the \textbf{Hardy space}, which is a subspace of the space 
$\mathcal{O}(U)$ of analytic functions
on a `certain domain' $U\subset \mathbb{C}$ which are defined by 
`certain growth conditions'. Roughly speaking,
`certain growth conditions' means that the boundary growth is controlled in 
an averaged $L^p$ sense. We usually consider
$U$ to be the unit disc $\D:=\{z\in\C: |z|<1\}$ or the upper half plane 
$\HH:=\{z\in\C: \im z>0\}$.

\begin{definition}
Let $0<p<\infty$. We define the \textbf{Hardy space over the unit disc $\D$} as:
\begin{equation*}
  H^p(\D):=\left\{ f\in \mathcal{O}(\D):\| f\|_{H^p}^p := \sup_{0<r<1}\frac{1}{2\pi}\int |f(re^{i\theta})|^p d
  \theta <\infty\right\}
\end{equation*}
If $p=\infty$, we define $H^\infty(\D)$ to be the space of bounded analytic functions on $\D$, i.e.
\begin{equation*}
  H^\infty(\D):=\{f\in \mathcal{O}(\D): \|f\|_{H^\infty}:=\sup_{z\in \D}|f(z)|<\infty\}
\end{equation*}
\end{definition}

\begin{definition}
Let $0<p<\infty$. We define the \textbf{Hardy space over the upper half plane $\HH$} as:
\begin{equation*}
  H^p(dt):=\left\{f\in \mathcal{O}(\HH): \|f\|_{H^p}^p:=\sup_{y>0}\int |f(x+iy)|^p dx <\infty\right\}
\end{equation*}
If $p=\infty$, we define $H^\infty(dt)$ to be the space of bounded analytic functions on $\HH$, i.e.
\begin{equation*}
  H^\infty(dt):=\{f\in \mathcal{O}(\HH): \|f\|_{H^\infty}:=\sup_{z\in \HH}|f(z)|<\infty\}
\end{equation*}
\end{definition}
These two definitions are closely related and many results holds true on $H^p(\D)$ will be also holds
true on $H^p(dt)$. In this thesis, we only focus on the formalization for $H^p(\D)$.

\begin{remark}
One can extend the definition of Hardy space over general (bounded) domain in $\mathbb{C}$ as follow:
\begin{definition}
Let $0<p<\infty$ and $U\subset \mathbb{C}$ be a (bounded) domain. We say an analytic function $f\in \mathcal{O}(U)$
over $U$ is a $H^p$ function if the function $|f|^p$ has a harmonic majorant, and we define $\|f\|^p_{H^p}$
to be the value of the least harmonic majorant at $z=0$. For $p=\infty$, $f\in \mathcal{O}(U)$ is
$H^\infty$ if it is a bounded analytic function over $U$ and $\|f\|_{H^\infty}:=\sup_{z\in U}|f(z)|$.
\end{definition}
Theorem \ref{Harmonic majorant criterion} shows that this is an equivalent definition with above definitions over $\D$.
\end{remark}


The Hardy space $H^p(\D)$ is a normed space for $p\geq 1$ with the
norm $\|\cdot\|_{H^p}$. Moreover, it is a metric space for all $0<p\leq \infty$,
with the metric $d(f,g)=\|f-g\|_{H^p}^p$ for $0<p<1$ and the metric induced by
the norm for $p\geq 1$. Moreover, we have:
\begin{theorem}
For $0<p\leq \infty$, $H^p$ is a complete metric space.
\end{theorem}
Similar as $L^p$ space, there is monotonicity for $H^p$ space as well:
\begin{theorem}
For $0<p\leq q\leq \infty$, we have $H^q\leq H^p$ as subspace.
\end{theorem}


\section{Blaschke product}
\begin{definition}
Given a sequence $\{z_n\}$ on $\mathbb{D}$, we say that this sequence satisfies
\textbf{Blaschke's condition} if $\sum (1-|z_n|)<\infty$. Given an analytic function $f$ on 
$\mathbb{D}$, we call this function $f$ satisfies \textbf{Blaschke's condition} if its
zeros $\{z_n\}$ satisfies Blaschke's condition.
\end{definition}

\begin{theorem}\label{Blaschke condition from harmonic majorant}
Let $f\not\equiv 0$ be an analytic function on $\mathbb{D}$ with zeros $\{z_n\}$. If
$\log|f|$ has a harmonic majorant, then $f$ satisfies Blaschke's condition:
\begin{equation*}
  \sum(1-|z_n|)<\infty.
\end{equation*}
If moreover $f(0)\neq 0$ and $u$ is the least harmonic majorant of $\log|f|$, then
\begin{equation*}
  \sum(1-|z_n|)\leq u(0)-\log|f(0)|.
\end{equation*}
\end{theorem}

\begin{theorem}\label{Blaschke product}
Let $\{z_n\}$ be a sequence in $\mathbb{D}$ satisfying Blaschke's condition, and let $m$ be the
number of $z_n$ equal to $0$. Then the \textbf{Blaschke product}
\begin{equation}\label{eq:Blaschke product}
  B(z)=z^m\prod_{|z_n|\neq 0}\left(-\frac{\overline{z_n}}{|z_n|}\right)
  \frac{z-z_n}{1-\overline{z_n}z}
\end{equation}
converges on $\mathbb{D}$. The function $B$ is in $H^\infty(\mathbb{D})$ with $|B(z)|\leq 1$,
and the zeros of $B$ are precisely the points $z_n$, each zero having multiplicity equal to
the number of times it occurs in $\{z_n\}$. Moreover,
\begin{equation*}
  |B(e^{i\theta})|=1
\end{equation*}
for almost every $\theta$, where $B(e^{i\theta})=\lim_{r\rightarrow 1}B(re^{i\theta})$ is the
radial limit of $B$.
\end{theorem}

Here's a characterization of Blaschke products in terms of harmonic majorants:
\begin{theorem}\label{characterization of Blaschke products}
Let $f\in H^\infty(\mathbb{D})$ with $\|f\|_{H^\infty}\leq 1$. The followings are equivalent:
\begin{itemize}
  \item[(a)] $f(z)=\lambda B(z)$ where $\lambda$ is constant with $|\lambda|=1$ and $B$ is a 
  Blaschke product;
  \item[(b)] $\lim_{r\rightarrow 1}\int \log |f(re^{i\theta})|d\theta/(2\pi)=0$;
  \item[(c)] The least harmonic majorant of $\log |f|$ is 0.
\end{itemize}
\end{theorem}

\section{Boundary functions of Hardy functions}
Every $H^p$ function has a boundary function in $L^p(\mathbb{T})$, and for $p\geq 1$ it can be
recovered from its boundary function via Poisson integral.
\begin{theorem}\label{thm:Hp_boundary_function}
Let $0<p\leq\infty$ and $f\in H^p(\D)$. Then $f$ has a nontangential limit $f^*(\zeta)$ at
almost every $\zeta\in\mathbb{T}$, and $f^*\in L^p(\mathbb{T}, \frac{d\theta}{2\pi})$ with
\begin{equation*}
  \|f^*\|_{L^p}=\|f\|_{H^p}.
\end{equation*}
Moreover:
\begin{itemize}
  \item[(a)] if $0<p<\infty$, then
  \begin{equation*}
    \lim_{r\rightarrow 1}\|f(re^{i\cdot})-f^*\|_{L^p}=0;
  \end{equation*}
  \item[(b)] if $1\leq p\leq\infty$, then $f=P[f^*]$.
\end{itemize}
\end{theorem}




\section{Nevanlinna class} 
\begin{definition}
An analytic function $f(z)$ on $\D$ is in the \textbf{Nevanlinna class} $N$ if 
\begin{equation*}
  \sup_r \int \log^+|f(r e^{i\theta})|\frac{d\theta}{2\pi}<\infty
\end{equation*}
\end{definition}
As $\log^+|f(z)|\leq (1/p)|f(z)|^p$, so $H^p\subset N$ for all $0<p\leq \infty$, i.e. Nevanlinna
class encompasses all Hardy spaces $H^p$. 

By Theorem \ref{Harmonic majorant criterion}, we can formulate the definition
in an equivalent definition:
\begin{definition}
An analytic function $f(z)$ on $\D$ is in the \textbf{Nevanlinna class} $N$ if 
the subharmonic function $\log^+|f(z)|$ has a harmonic majorant.
\end{definition}

\begin{remark}\label{rem:N_algebra}
Compared to $H^p(\D) \subset \mathcal{O}(\D)$,
which is just a subspace of $\mathcal{O}(\D)$ for $p<\infty$, $N\subset \mathcal{O}(\D)$ is
also a subalgebra of $\mathcal{O}(\D)$: if $f,g\in N$, then $fg\in N$ as
\begin{equation*}
  \log^+|fg|\leq \log^+|f|+\log^+|g|.
\end{equation*}
Moreover, if $f\in N$ has no zeros in $\D$, then $1/f\in N$. Indeed, let $U$ be a harmonic
majorant of $\log^+|f|$. Since $f$ has no zeros, $\log|f|$ is harmonic on $\D$, and
\begin{equation*}
  \log^+|1/f|=\log^+|f|-\log|f|\leq U-\log|f|,
\end{equation*}
so $\log^+|1/f|$ has the harmonic majorant $U-\log|f|$.
The closure of multiplication and inversion will sort of provide a more natural place for us to 
describe `factorization problem'.
\end{remark}

\begin{theorem}\label{characterization of Nevanlinna using log |f|}
Let $f(z)$ be an analytic function on $\mathbb{D}$ s.t. $f\not\equiv 0$. 
Then $f\in N$ if and only if $\log|f(z)|$ has least harmonic majorant
the Poisson integral of some finite measure on $\mathbb{T}$.
\end{theorem}
\begin{proof}

\end{proof}


\begin{lemma}\label{least harmonic majorant of log|f| for Nevanlinna f}
Let $f\in N$ and $f\not\equiv 0$. Let $B$ be the Blaschke product formed from zeros of $f$. Then:
\begin{itemize}
  \item $B$ converges;
  \item $g=f/B\in N$;
  \item $\log |g|$ is the least harmonic majorant of $\log |f|$.
\end{itemize}
\end{lemma}
\begin{proof}[Sketch of proof]
As $f\in N$, $\log |f|$ has a harmonic majorant, thus by Theorem \ref{Blaschke condition from harmonic majorant},
$f$ satisfies Blaschke's condition, thus $B$ converges. Now let $u$ be the least harmonic majorant
of $\log |f|$. Recall in Example \ref{log |f| is subharmonic}, $g$ as a zero-free analytic function
over $\mathbb{D}$, $\log |g|$ is harmonic on $\mathbb{D}$. As $|B(z)|\leq 1$, we have 
$\log |f(z)|\leq \log |g(z)|$, thus 
\begin{equation*}
  u(z)\leq \log|g(z)|.
\end{equation*}
On the other hand, 
\begin{equation*}
  \log |B(z)|=\log |f(z)|-\log|g(z)|\leq u(z)-\log |g(z)|
\end{equation*}
Recall Theorem \ref{characterization of Blaschke products}, the least harmonic majorant of $\log|B|$
is 0, thus 
\begin{equation*}
  0\leq u(z)-\log|g(z)|
\end{equation*}
which implies $u(z)=\log |g(z)|$. It's easy to check that $g\in N$ by definition.
\end{proof}


\begin{theorem}\label{nontangential limit of Nevanlinna}
Let $f\in N$ and $f\not\equiv 0$. Then $f$ has a nontangential limit $f^*(\zeta)$ at almost
every $\zeta\in\mathbb{T}$, and $\log|f^*|\in L^1(\mathbb{T})$.

Moreover, the least harmonic majorant of $\log|f|$ is $P[\mu]$, where
\begin{equation*}
  d\mu=\log|f^*|\,\frac{d\theta}{2\pi}+d\mu_s,
  \qquad \mu_s\perp\frac{d\theta}{2\pi}.
\end{equation*}
\end{theorem}

\subsection{Outer functions and Singular functions}
Before we discuss the main topic of this section, which is the factorization problem of
Nevanlinna functions, we'd like to first introduce two classes of functions: 
\textbf{outer functions} and \textbf{singular functions}.
\begin{definition}\label{def:outer_function}
An analytic function $F$ on $\D$ is an \textbf{outer function} if
\begin{equation*}
  F(z)=C\exp\left(\int K_z(\theta)\,l(\theta)\,\frac{d\theta}{2\pi}\right),
  \qquad z\in\D,
\end{equation*}
for some real-valued $l\in L^1(\mathbb{T})$ and some constant $C\in\mathbb{C}$ with $|C|=1$.
Here $K_z(\theta)=\frac{e^{i\theta}+z}{e^{i\theta}-z}$ is the Herglotz-Riesz kernel.
\end{definition}

\begin{theorem}\label{thm:outer_function_properties}
Let $F$ be an outer function on $\mathbb{D}$. Then:
\begin{itemize}
  \item[(a)] $F$ is analytic on $\D$ and has no zeros in $\D$;
  \item[(b)] $\log|F|=P[l]$ on $\D$;
  \item[(c)] $F\in N$;
  \item[(d)] $F$ has a nontangential limit $F^*(\zeta)$ at almost every $\zeta\in\mathbb{T}$, and
  $|F^*|=e^{l}$ almost everywhere on $\mathbb{T}$;
  \item[(e)] if $F_1,F_2$ are outer functions with $|F_1^*|=|F_2^*|$ almost everywhere on
  $\mathbb{T}$, then $F_1=cF_2$ for some constant $c\in\mathbb{C}$ with $|c|=1$.
\end{itemize}
\end{theorem}
\begin{proof}
Write $F=Ce^{G}$ with
\begin{equation*}
  G(z)=\int K_z(\theta)\,l(\theta)\,\frac{d\theta}{2\pi}.
\end{equation*}

(a) For $|z|\leq r<1$ we have $|K_z(\theta)|\leq\frac{1+r}{1-r}$, and $K_z(\theta)$ is analytic
in $z$. Hence $G$ is analytic on $\D$ and so is
$F=Ce^{G}$. Since the exponential has no zeros, neither does $F$.

(b) As $|C|=1$, $l$ is real-valued and $\text{Re } K_z(\theta)=P_z(\theta)$, note
that $\log |e^G|=\text{Re } G$,

\begin{equation*}
  \log|F(z)|=\text{Re } G(z)=\int P_z(\theta)\,l(\theta)\,\frac{d\theta}{2\pi}=P[l](z).
\end{equation*}

(c) Write $l^+=\max(l,0)$. By Lemma \ref{Properties of Poisson kernel}(a), $P_z>0$, so
$\log|F|=P[l]\leq P[l^+]$, and since $P[l^+]\geq 0$, also $\log^+|F|\leq P[l^+]$. Using
$P_{re^{i\theta}}(\varphi)=P_{re^{i\varphi}}(\theta)$, Fubini's theorem and Lemma
\ref{Properties of Poisson kernel}(a) again,
\begin{equation*}
  \frac{1}{2\pi}\int_0^{2\pi}\log^+|F(re^{i\theta})|\,d\theta
  \leq\frac{1}{2\pi}\int_0^{2\pi}P[l^+](re^{i\theta})\,d\theta
  =\frac{1}{2\pi}\int_0^{2\pi}l^+(\varphi)\,d\varphi
  \leq\|l\|_{L^1}
\end{equation*}
for every $0<r<1$. Hence $F\in N$.

(d) By (c) and $F\not\equiv0$, Theorem \ref{nontangential limit of Nevanlinna} shows that $F$ has
a nontangential limit $F^*$ almost everywhere. By (b) and Corollary
\ref{cor:poisson_nontangential_ae}, $\log|F|=P[l]$ has nontangential limit $l$ almost
everywhere. At almost every $\zeta\in\mathbb{T}$ both limits exist, and then
\begin{equation*}
  |F^*(\zeta)|=\lim_{\substack{z\to\zeta\\ z\in\Gamma_\alpha(\zeta)}}e^{\log|F(z)|}=e^{l(\zeta)}.
\end{equation*}

(e) Let $F_j=C_je^{G_j}$, where $G_j$ is defined from $l_j$ as above, $j=1,2$. By (d),
$e^{l_1}=|F_1^*|=|F_2^*|=e^{l_2}$ almost everywhere, so $l_1=l_2$ almost everywhere. Hence
$G_1=G_2$ on $\D$, and $F_1=cF_2$ with $c=C_1/C_2$, $|c|=1$.
\end{proof}

\begin{remark}
One can construct an outer function as follow: given $h:\mathbb{T}\rightarrow [0,\infty)$
s.t. $\log h\in L^1(\mathbb{T})$
(think of it as the modulos of the boundary function of Nevanlinna functions over $\mathbb{D}$, 
see \ref{nontangential limit of Nevanlinna}), then set $u(z)=P[\log h](z)$ and $v(z)$ to be one of the
harmonic conjugates of $u(z)$. Then the function $F(z):= e^{u(z)+i v(z)}$ is an outer function,
and it is determined by $h(\theta)$ except for a unimodular constant factor due to the choice
of the conjugate $v(z)$. 
\end{remark}

\begin{definition}
An analytic function $S$ on $\D$ is a \textbf{singular function} if
\begin{equation*}
  S(z)=\exp\left(-\int K_z(\theta)\,d\mu(\theta)\right),
  \qquad z\in\D,
\end{equation*}
for some finite singular positive measure $\mu$ on $\mathbb{T}$.
\end{definition}

\begin{lemma}\label{lem:log_singular_function}
For singular function $S(z)=\exp\left(-\int K_z(\theta)\,d\mu(\theta)\right)$
 on $\mathbb{D}$, we have $\log|S|=-P[\mu]$ on $\mathbb{D}$.
\end{lemma}
\begin{proof}
Given $z\in\D$, as $|K_z(\theta)|\leq\frac{1+|z|}{1-|z|}$ and $\mu$ is finite, 
$\int K_z(\theta)\,d\mu(\theta)$ converges. As $\log|e^w|=\text{Re } w$, $\mu$ is a positive
measure and $\text{Re } K_z(\theta)=P_z(\theta)$,
\begin{equation*}
  \log|S(z)|=-\text{Re }\int K_z(\theta)\,d\mu(\theta)
  =-\int P_z(\theta)\,d\mu(\theta)=-P[\mu](z).
\end{equation*}
\end{proof}

\begin{lemma}\label{lem:singular_function_in_N}
If $S$ is a singular function on $\mathbb{D}$, then $S\in N$.
\end{lemma}
\begin{proof}
By Lemma \ref{lem:log_singular_function} and Lemma \ref{Properties of Poisson kernel}(a),
$\log|S|=-P[\mu]\leq 0$ on $\D$, as $\mu\geq 0$ and $P_z>0$. Hence $\log^+|S|=0$ on $\D$,
so $S\in N$.
\end{proof}

\begin{theorem}\label{equivalent characterization of singular function}
Let $S$ be an analytic function on $\D$. Then $S$ is a singular function if and only if $S$ has the
following properties:
\begin{itemize}
  \item[(a)] $S$ has no zeros in $\D$;
  \item[(b)] $|S(z)|\leq 1$ for all $z\in\D$;
  \item[(c)] $S$ has a nontangential limit $S^*(\zeta)$ at almost every $\zeta\in\mathbb{T}$, and
  $|S^*|=1$ almost everywhere on $\mathbb{T}$;
  \item[(d)] $S(0)>0$.
\end{itemize}
\end{theorem}

\begin{proof}[Proof of one direction]
($\Rightarrow$) Let $S(z)=\exp\left(-\int K_z(\theta)\,d\mu(\theta)\right)$ with $\mu\geq 0$
finite and singular.
(a) holds as exponential has no zeros.
(b) was shown in the proof of Lemma \ref{lem:singular_function_in_N}.
(c) By Lemma \ref{lem:singular_function_in_N} and Theorem \ref{nontangential limit of Nevanlinna},
$S$ has a nontangential limit $S^*$ almost everywhere. By Lemma \ref{lem:log_singular_function}
and Theorem \ref{thm:singular_poisson_nontangential_zero}, $\log|S|=-P[\mu]$ has nontangential
limit $0$ almost everywhere. So at a.e. $\zeta\in\mathbb{T}$,
\begin{equation*}
  |S^*(\zeta)|=\lim_{\substack{z\to\zeta\\ z\in\Gamma_\alpha(\zeta)}}e^{-P[\mu](z)}=1.
\end{equation*}
(d) holds because $K_0(\theta)=1$ gives $S(0)=e^{-\mu(\mathbb{T})}>0$.

We will skip the proof of ($\Leftarrow$) as it's not needed in the project.
\end{proof}



\subsection{Canonical factorization theorem of Nevanlinna funcitons}
Now we'd like to discuss the \textit{factorization problem} of Nevanlinna functions.
Given $f\in N$ with $f\not\equiv 0$, consider $B$: Blaschke product formed by
its zeros and $g=f/B$. By Lemma \ref{least harmonic majorant of log|f| for Nevanlinna f},
we know that $\log|g|$ is the least harmonic majorant of $\log |f|$. Recall 
Theorem \ref{characterization of Nevanlinna using log |f|}, we can write 
$\log |g|=P[\mu]$ for some finite measure $\mu$. By Theorem \ref{nontangential limit of Nevanlinna},
we can write:
\begin{equation*}
  d\mu=\log|f^*|\,\frac{d\theta}{2\pi}+d\mu_s,
  \qquad \mu_s\perp\frac{d\theta}{2\pi}.
\end{equation*}
Now take Jordan decomposition on singular measure
 $\mu_s$, we get the unique decomposition:
 \begin{equation*}
  \mu_s=-\mu_1+\mu_2
 \end{equation*}
 where $\mu_1$ and $\mu_2$ are positive singular measure that are mutually singular to each other. Thus we 
write:
\begin{equation*}
  d\mu=\log|f^*|\,\frac{d\theta}{2\pi}-d\mu_1+d\mu_2.
\end{equation*}

Now consider the \textbf{outer function}
\begin{equation*}
  \tilde{F}(z)=\exp\left(\int K_z(\theta)\log |f^*(e^{i\theta})|\frac{d\theta}{2\pi}\right)
\end{equation*}
Here $K_z(\theta)=\frac{e^{i\theta}+z}{e^{i\theta}-z}$ is the Herglotz-Riesz kernel. Moreover, consider
the \textbf{singular functions}:
\begin{equation*}
  S_j(z)=\exp\left(-\int K_z(\theta)d\mu_j(\theta)\right),\quad j=1,2.
\end{equation*}
Recall that $\text{Re } K_z(\theta)=P_z(\theta)$, which is the Poisson kernel, and
$\log |e^w|=\text{Re } w$, so 
\begin{equation*}
  \log |\tilde{F}(z)|=P[\log |f|](z),\; \log |S_j(z)|=- P[\mu_j](z),\quad j=1,2.
\end{equation*}
Therefore, 
\begin{align*}
  \log |g(z)|=P[\mu](z)&=P[\log |f|](z)+P[\mu_1](z)-P[\mu_2](z)\\
  &= \log |\tilde{F}(z)|+\log |S_1(z)|-\log|S_2(z)|.
\end{align*}
Then we have:
\begin{equation*}
  \log \left|\frac{g(z) S_2(z)}{\tilde{F}(z)S_1(z)}\right|=0
\end{equation*}
which takes $\exp$ on both sides yields:
\begin{equation*}
  \left|\frac{f(z)S_2(z)}{B(z)\tilde{F}(z)S_1(z)}\right|=1,\quad \forall z\in \mathbb{D}.
\end{equation*}
By maximum modulus principle, 
\begin{equation*}
  \frac{f(z)S_2(z)}{B(z)\tilde{F}(z)S_1(z)}=C,\quad C\in\mathbb{C},\; |C|=1.
\end{equation*}
Set $F=C\tilde{F}$. By Definition \ref{def:outer_function}, $F$ is again an outer function, and
$f=BFS_1/S_2$.
This gives most of the \textbf{canonical factorization theorem}:
\begin{theorem}
Let $f\in N$, $f\not\equiv 0$. Then
\begin{equation*}
  f(z)=B(z)F(z)S_1(z)/S_2(z),
\end{equation*}
where $B$ is Blaschke product formed by zeros of $f$,
$F$ is an outer function, and $S_1,S_2$ are
singular functions whose positive singular measures $\mu_1,\mu_2$ are mutually singular.
This factorization is unique.
Conversely, every function of this form is in $N$.
\end{theorem}
\begin{proof}
The existence part is given as above.

For uniqueness, the Blaschke product $B$ is determined by the zeros of $f$. Suppose
$f=BFS_1/S_2=BF'S_1'/S_2'$, where $S_1,S_2$ come from $\mu_1,\mu_2$ and $S_1',S_2'$ from
$\nu_1,\nu_2$. Since the zeros of $B\not\equiv 0$ are isolated, $FS_1/S_2=F'S_1'/S_2'$ on $\D$
minus a discrete set, hence on $\D$ by continuity. By Theorem
\ref{equivalent characterization of singular function}(c), $|S_1^*|=|S_2^*|=|S_1'^*|=|S_2'^*|=1$
almost everywhere, so $|F^*|=|F'^*|$ almost everywhere, and by Theorem
\ref{thm:outer_function_properties}(e), $F=cF'$ for a constant $c$ with $|c|=1$. Thus
$cS_1/S_2=S_1'/S_2'$, and taking $\log|\cdot|$, Lemma \ref{lem:log_singular_function} gives
\begin{equation*}
  P[\mu_2-\mu_1]=P[\nu_2-\nu_1]\quad\text{on }\D.
\end{equation*}
Since a finite measure on $\mathbb{T}$ is determined by its Poisson integral,
$\mu_2-\mu_1=\nu_2-\nu_1$. Both $(\mu_2,\mu_1)$ and $(\nu_2,\nu_1)$ are pairs of mutually
singular positive measures, so by uniqueness of the Jordan decomposition, $\mu_1=\nu_1$ and
$\mu_2=\nu_2$. Hence $S_j=S_j'$ for $j=1,2$, so $c=1$ and $F=F'$.

For converse part Since $N$ is closed under multiplication by Remark \ref{rem:N_algebra}, it suffices to show that
$B,F,S_1,1/S_2\in N$. We have $B\in H^\infty\subset N$ by Theorem \ref{Blaschke product},
$F\in N$ by Theorem \ref{thm:outer_function_properties}(c), and $S_1\in N$ by Lemma
\ref{lem:singular_function_in_N}. Likewise $S_2\in N$, and $S_2$ has no zeros in $\D$, so
$1/S_2\in N$ by Remark \ref{rem:N_algebra}.
\end{proof}
